% =============================================================
% Sheaf Bandwidth: A Categorical Bridge Between Topological
% Signal Processing and Finite Element Exterior Calculus
%
% Paper v0.1 (LaTeX source)
% =============================================================

\documentclass[11pt,reqno]{amsart}

% ---------------- Geometry & spacing ----------------
\usepackage[margin=1in]{geometry}
\usepackage{microtype}

% ---------------- Math packages ----------------
\usepackage{amsmath}
\usepackage{amssymb}
\usepackage{amsthm}
\usepackage{mathtools}
\usepackage{mathrsfs}     % \mathscr for sheaves
\usepackage{bbm}          % \mathbbm for blackboard with non-letters

% ---------------- Diagrams (optional, for sheaf restriction maps) ----------------
\usepackage{tikz-cd}

% ---------------- Lists & tables ----------------
\usepackage{enumitem}
\usepackage{booktabs}

% ---------------- Hyperref (load last) ----------------
\usepackage[colorlinks=true,linkcolor=blue!50!black,citecolor=blue!50!black,urlcolor=blue!50!black]{hyperref}
\usepackage{cleveref}

% =============================================================
% Theorem environments
% =============================================================
\theoremstyle{plain}
\newtheorem{theorem}{Theorem}[section]
\newtheorem{lemma}[theorem]{Lemma}
\newtheorem{proposition}[theorem]{Proposition}
\newtheorem{corollary}[theorem]{Corollary}

\theoremstyle{definition}
\newtheorem{definition}[theorem]{Definition}
\newtheorem{example}[theorem]{Example}
\newtheorem{remark}[theorem]{Remark}

% Empirical Proposition: a non-theorem-strength statement supported by
% (i) a heuristic skeleton and (ii) crossed evidence channels.
\theoremstyle{definition}
\newtheorem{empirical}[theorem]{Empirical Proposition}

% Open Problem: an explicit, independent open question.
\theoremstyle{definition}
\newtheorem{openproblem}[theorem]{Open Problem}

% Cleveref naming
\crefname{theorem}{Theorem}{Theorems}
\crefname{lemma}{Lemma}{Lemmas}
\crefname{proposition}{Proposition}{Propositions}
\crefname{corollary}{Corollary}{Corollaries}
\crefname{definition}{Definition}{Definitions}
\crefname{remark}{Remark}{Remarks}
\crefname{empirical}{Empirical Proposition}{Empirical Propositions}
\crefname{openproblem}{Open Problem}{Open Problems}

% =============================================================
% Symbol macros (sync with paper/v0.1/notes/symbols.md §4)
% =============================================================

% --- Operators ---
\DeclareMathOperator{\sheafBW}{sheafBW}
\DeclareMathOperator{\spec}{spec}
\DeclareMathOperator{\rank}{rank}
\DeclareMathOperator{\im}{im}
\DeclareMathOperator{\eval}{eval}

% --- Subscript / superscript labels ---
\newcommand{\samp}{\mathrm{samp}}
\newcommand{\fin}{\mathrm{fin}}
\newcommand{\TopSP}{\mathrm{TopSP}}
\newcommand{\FEM}{\mathrm{FEM}}
\newcommand{\BZ}{\mathrm{BZ}}
\newcommand{\occ}{\mathrm{occ}}

% --- Categories ---
\newcommand{\Hilb}{\mathsf{Hilb}}
\newcommand{\Vect}{\mathsf{Vect}}
\newcommand{\CellSh}{\mathsf{Cell\text{-}Sh}}

% --- Number sets ---
\newcommand{\C}{\mathbb{C}}
\newcommand{\R}{\mathbb{R}}
\newcommand{\N}{\mathbb{N}}
\newcommand{\Z}{\mathbb{Z}}

% --- Sampling sheaf shorthand ---
\newcommand{\samps}{\mathscr{A}^{\samp}}

% =============================================================
% Lean 4 formalization cross-reference
% =============================================================
% \leanref{file}{name} — footnote pointing to Lean 4 formalization.
% Example: \leanref{MRI/Cellular/SamplingSpectral}{strict\_equality}
% Companion artifact: lean/CROSSREF.md (full paper-Lean correspondence table).
\newcommand{\leanref}[2]{%
  \footnote{Formalized in Lean 4: theorem \texttt{#2} in file
    \texttt{lean/#1.lean}. Full correspondence: \texttt{lean/CROSSREF.md}.}}

% =============================================================
% Title & metadata
% =============================================================
\title[Sheaf Bandwidth: A Categorical Bridge]{Sheaf Bandwidth:\\
  A Categorical Bridge Between Topological Signal Processing\\
  and Finite Element Exterior Calculus}

\author{[Author placeholder]}
\address{Independent Researcher}
\email{[Email placeholder]}

\date{\today}

\subjclass[2020]{18F20, 65N30, 94A20, 55U10}
\keywords{Sheaf bandwidth, cellular sheaf, finite element exterior calculus,
  Hodge decomposition, sampling theorem, Brillouin zone tetrahedron,
  $k$-space sampling, MRI}

% =============================================================
\begin{document}
\maketitle

% Phase 1 sanity: \nocite{*} forces bibtex to include all entries
% so we exercise every refs.bib entry. Remove before submission.
\nocite{*}

% =============================================================
\begin{abstract}
% =============================================================
% Abstract (120-180 words, five sentences)
% Template: paper/v0.1/notes/narrative.md §2
% =============================================================

A cellular sheaf of Hilbert spaces equips a regular cell complex with
stalks and restriction maps, and Robinson \cite[\S 6]{Robinson2014}
introduced an analytic invariant of such sheaves---the
\emph{bandwidth}---while leaving open whether this invariant admits a
definition intrinsic to the sheaf rather than to a chosen ambient
representation.
We propose
$\sheafBW(\mathscr{A}) := \sqrt{\gamma_+(\delta^0)/\Gamma_+(\delta^0)}$,
prove unitary invariance (\cref{prop:sheafBW-unitary}), and establish
a strict equality
$\sheafBW(\samps) = \sigma_{\min}(A)/\sigma_{\max}(A)$
(\cref{thm:strict-equality}) on an abstract sampling cell complex
$X^A$ with a canonical double $0$-cell construction. The framework
is instantiated in three previously disjoint domains---topological
signal processing on simplicial complexes, finite element
electromagnetics, and Brillouin zone tetrahedron quadrature in
condensed matter---unified by a common categorical template, with a
triangle equivalence (\cref{thm:triangle-equiv}) connecting finite
element exterior calculus singular value bounds, topological signal
processing Riesz frame conditions, and Robinson sheaf cohomology;
numerical verification across three toy operators matching the common
finite-dimensional template confirms the strict equality to within the
expected floating-point error budget ($\mathrm{reldiff} \le 1.64
\times 10^{-14}$, driven by the squaring of the condition number under
the $A^* A$ form). An empirical
proposition (\cref{empirical:8b}) relates the sampling and
integration maps via a Pareto trade-off, supported by two
empirical channels (a project-internal numerical study and five
decades of Brillouin zone practice) and a Fourier-spectrum
heuristic. Two independent open problems remain: a
non-trivial dagger structure on the full category
(\cref{op:non-trivial-dagger}) and a constructive Pareto proof of the
empirical proposition (\cref{op:8b-rigor}).

\end{abstract}

% =============================================================
% Section 1. Introduction
% =============================================================
\section{Introduction}\label{sec:intro}
% =============================================================
% Section 1. Introduction
% Template: paper/v0.1/notes/narrative.md §4 §1
% Structure: §1.1 Robinson open problem / §1.2 baseline black-box
%            §1.3 FEEC × sheaf SP disconnect / §1.4 contributions / §1.5 outline
% =============================================================

% -------------------------------------------------------------
\subsection{Robinson's open problem on sheaf bandwidth}\label{subsec:intro-robinson}

Cellular sheaves of Hilbert spaces, as developed by
\cite{Robinson2014} and refined in the spectral theory of
\cite{HansenGhrist2019}, equip a regular cell complex $X$ with stalks
$\mathscr{A}(\sigma)$ on each cell and bounded restriction maps along
incidences. Robinson's monograph closes \cite[\S 6]{Robinson2014} with
an explicit open question: does sheaf bandwidth admit a definition
that is intrinsic to the sheaf, not merely to a chosen ambient
representation? In our reading, the question splits into three sub-questions:
\begin{enumerate}[leftmargin=*,label=(Q\arabic*)]
\item \textbf{Categorical intrinsicness.} Is there a real-valued
  invariant $\sheafBW(\mathscr{A})$, defined from the sheaf data alone
  (stalks, restriction maps, sheaf Laplacian), that is preserved under
  unitary equivalence in the ambient cellular sheaf category?
\item \textbf{Spectral handle.} Does this invariant admit a closed
  form on the spectrum of the sheaf Laplacian
  $L^0 = (\delta^0)^* \delta^0$, comparable in tractability to the
  singular value ratio $\sigma_{\min}/\sigma_{\max}$ of an ambient
  matrix?
\item \textbf{Concrete reduction.} On a canonical class of sheaves
  arising from a sampling problem, does the intrinsic invariant
  reduce \emph{strictly} (not just up to constants) to a
  domain-specific singular value ratio?
\end{enumerate}
This paper answers all three affirmatively, in the
finite-dimensional, spectral-gap, unitary-invariant version of the
question as split here
(\cref{def:sheafBW,prop:sheafBW-unitary,thm:strict-equality}), and
exhibits the answer in three independent application domains
(\cref{sec:three_instances}). The technical formulation in
\cref{subsec:tri-narrowing} narrows these three sub-questions to a
working two-clause specification, in which (\textup{Q2}) is absorbed
into the spectral-gap form of \cref{def:sheafBW} itself. The
infinite-dimensional extension and any categorical generalization
beyond the present cellular-sheaf-of-finite-dimensional-Hilbert-spaces
setting are explicitly out of scope; see \cref{subsec:tri-edge} for
the four boundary conditions (L1)--(L4).

% -------------------------------------------------------------
\subsection{The black-box treatment of $k$-space sampling in MRI}\label{subsec:intro-mri-blackbox}

Magnetic resonance imaging reconstruction by 2026 rests on four
schools that treat the choice of sampling geometry as a black box
upstream of reconstruction:
\begin{itemize}[leftmargin=*]
\item \emph{Compressed sensing} \cite{LustigDonohoPauly2007} treats
  sampling as a fixed acquisition pattern and optimizes only the
  reconstruction prior;
\item \emph{Topological signal processing}
  \cite{BarbarossaSardellitti2020} treats sampling on simplicial
  complexes via Riesz frame inequalities, again with the geometry
  given;
\item \emph{Sparkling-style point process design}
  \cite{LazarusEtAl2019Sparkling} treats sampling design as a
  point-process optimization detached from the reconstruction
  operator;
\item \emph{Sheaf signal processing}
  \cite{Robinson2014,HansenGhrist2019,HansenGhrist2019Allerton}
  studies sheaf bandwidth as an analytic invariant without exposing a
  computable reduction to acquisition-side singular values.
\end{itemize}
Each school is internally precise. None exposes a single scalar
that is at once (i) intrinsic to the abstract sampling structure and
(ii) directly computable from the acquisition matrix. We will refer
to this gap as the \emph{sampling black box} and resolve it via the
triangle equivalence of \cref{thm:triangle-equiv}.

% -------------------------------------------------------------
\subsection{The FEEC and cellular-sheaf disconnect}\label{subsec:intro-disconnect}

Finite element exterior calculus (FEEC), as systematized by
\cite{AFW2010,FalkWinther2014} on the algebraic side and by
\cite{BochevHyman2006,Christiansen2010} on the construction side,
provides closed Hilbert complex theory, Hodge decomposition with
stalk-wise harmonic representatives, and bounded cochain projections
of arbitrary commuting type. Cellular sheaf signal processing,
beginning with \cite{Robinson2014}, develops cohomology and bandwidth
on the same combinatorial substrate. Despite the shared substrate of
cell complexes equipped with cochain spaces, the two literatures have
not been merged in print: \cite{ChristiansenRapetti2025} surveys 30
years of FEEC without referencing sheaf signal processing, and
\cite{CerejeirasKahlerLegatiuk2025} surveys recent topological signal
processing without invoking AFW Hilbert complexes. We attribute this
to a missing concrete bridge---a single object on which the two
machineries meet and produce the same scalar---which we supply via
the abstract sampling cell complex $X^A$ of \cref{def:Xsamp}.

% -------------------------------------------------------------
\subsection{Contributions}\label{subsec:intro-contributions}

The contributions of this paper are the following.

\begin{enumerate}[leftmargin=*,label=(C\arabic*)]
\item \textbf{Triangle Equivalence Theorem
    (\cref{thm:triangle-equiv}).} Three formulations of well-conditioning
  for $k$-space sampling---FEEC singular value bounds, topological
  signal processing Riesz frame conditions, and Robinson sheaf
  cohomology with sheaf bandwidth---are formally equivalent
  \emph{under the common finite-dimensional sampling-operator normal
  form} of \cref{subsec:tri-setup}; the equivalence is not claimed
  to hold in the original infinite-dimensional or qualitative
  settings of the three source literatures. The proof of the
  equivalence is decomposed into a linear-algebraic part
  $(a) \Leftrightarrow (b)$ (\cref{subsec:tri-ab}) and a
  sheaf-cohomological part $(a) \Leftrightarrow (c)$
  (\cref{subsec:tri-ac}).

\item \textbf{Sheaf bandwidth, categorically intrinsic
    (\cref{def:sheafBW,prop:sheafBW-unitary}).} The proposed invariant
  $\sheafBW(\mathscr{A}) := \sqrt{\gamma_+(\delta^0)/\Gamma_+(\delta^0)}$
  uses only the sheaf Laplacian spectral gap and is unitarily
  invariant in the cellular sheaf category. This addresses
  sub-questions (\textup{Q1}) and (\textup{Q2}) of
  \cref{subsec:intro-robinson}.

\item \textbf{Strict equality on the abstract sampling cell complex
    (\cref{thm:strict-equality}).} On the canonical sampling sheaf
  $\samps$ over the cell complex $X^A$
  (\cref{def:Xsamp,def:sampling-sheaf}), constructed with a double
  $0$-cell $\{\sigma_B, \sigma_\infty\}$ to satisfy the regularity
  conditions of \cite[Def.~1, Cond.~4]{HansenGhrist2019}, the
  intrinsic invariant reduces strictly to
  $\sigma_{\min}(A)/\sigma_{\max}(A)$ on the well-conditioned locus
  $\sigma_{\min}(A) > 0$ (equivalently $H^0(\samps) = 0$); the
  rank-deficient boundary is recorded in
  \cref{rem:rank-deficient-locus}. This addresses sub-question
  (\textup{Q3}) of \cref{subsec:intro-robinson}.

\item \textbf{Three-instance unification by a common categorical
    template (\cref{sec:three_instances}).} Three previously disjoint
  domains fit the same template
  $V_? \subset \C^N$ + $A : V_? \to \C^M$:
  topological signal processing
  ($V_?  = V_B^{\TopSP}$, the bandlimited subspace),
  finite element electromagnetics
  ($V_?  = V_S^{\FEM}$, a finite-dimensional restriction of an AFW
  conforming subspace), and Brillouin zone tetrahedron quadrature
  ($V_?  = V_{\occ}^{\BZ}$, the Whitney $0$-form span over occupied
  $k$-points; see \cref{subsec:6-bz} for scope).
  \Cref{thm:strict-equality} applies verbatim in each.

\item \textbf{Empirical proposition on sampling--integration
    trade-off (\cref{empirical:8b}).} The sampling map $S$
  (\cref{def:sampling-map}) and the integration map $I$
  (\cref{def:integration-map}) trace a Pareto frontier under
  geometry optimization: no admissible $V$ simultaneously maximizes
  $S$ and minimizes $I$. The supporting evidence comprises two
  empirical channels and a heuristic skeleton: a project-internal
  numerical study on $35$ MRI sampling configurations with a
  multi-$N$ multi-$K$ robustness stratification (\cref{sec:numerical});
  five decades of Brillouin zone practice in condensed matter
  \cite{Blochl1994,LehmannTaut1972,JepsenAndersen1971,MonkhorstPack1976},
  on the integration side; and the Euclidean Fourier-spectrum
  framework of \cite{Pilleboue2015}, contributing directional
  content only rather than direct evidence on cellular-sheaf
  integration.

\item \textbf{Numerical verification of the strict equality
    (\cref{sec:numerical}).} The strict equality of
  \cref{thm:strict-equality} is verified across three toy operators
  matching the common finite-dimensional template, with maximum
  relative deviation $1.64 \times 10^{-14}$ (Instance~2; Instances~1
  and~3 are within $8 \times 10^{-16}$). The deviation reflects the
  expected roundoff amplification of comparing
  $\sqrt{\lambda_{\min}(A^* A)/\lambda_{\max}(A^* A)}$ with
  $\sigma_{\min}(A)/\sigma_{\max}(A)$ on conditioned matrices, since
  $\kappa(A^* A) = \kappa(A)^2$. The verification is honest in scope:
  it confirms the linear-algebraic identity and the code path, not the
  underlying physical evidence per instance, which is recorded
  separately in \cref{tab:7-evidence}.

\item \textbf{Two independent open problems
    (\cref{op:non-trivial-dagger,op:8b-rigor}).} The categorical
  bridge constructed here leaves two precisely formulated open
  questions: a non-trivial dagger structure on the full category
  $\CellSh^{\fin}_{\Hilb}(X)$ extending the unitary-subcategory
  groupoid (\cref{prop:unitary-groupoid}); and a constructive
  Pareto-optimality proof upgrading \cref{empirical:8b} to a
  theorem via compactness, convex relaxation, and spectral
  duality.
\end{enumerate}

% -------------------------------------------------------------
\subsection{Paper outline}\label{subsec:intro-outline}

\Cref{sec:related} reviews the four schools whose tools we combine:
finite element exterior calculus, cellular sheaf signal processing,
topological signal processing, and dagger Galerkin theory.
\Cref{sec:triangle} states and proves the triangle equivalence theorem
and produces a candidate sheaf bandwidth that is not yet categorically
intrinsic. \Cref{sec:feec_review} reviews the FEEC machinery
(closed Hilbert complex, Hodge decomposition, bounded cochain
projection) needed to lift bandwidth to a categorically intrinsic
invariant, and presents a four-school comparison table.
\Cref{sec:sheafbw} constructs the categorically intrinsic candidate
$\sheafBW$, proves unitary invariance, builds the abstract sampling
cell complex $X^A$ with its double $0$-cell, and establishes the
strict equality \cref{thm:strict-equality}.
\Cref{sec:three_instances} instantiates the framework in topological
signal processing, finite element electromagnetics, and Brillouin
zone quadrature, states the empirical proposition on the
$S$--$I$ Pareto trade-off, and records the constructive Pareto
proof as an independent open problem. \Cref{sec:numerical} verifies
the strict equality numerically and states the scope of the
verification. \Cref{sec:discussion} restates the seven contributions,
discusses the two open problems, and outlines the engineering
verification on the MRI side that lies outside this paper.

\medskip\noindent
The three motivating threads---Robinson's open question, the MRI
black-box treatment of sampling, and the FEEC and cellular-sheaf
disconnect---are addressed jointly by a single object on which all
three meet. \Cref{sec:related} reviews the four schools whose
tools we use, and \cref{sec:triangle} begins the construction with a
triangle equivalence theorem.

\paragraph{Formalization in Lean 4.}
The main finite-dimensional operator/sheaf-core theorems, propositions,
and corollaries of \cref{sec:triangle} and \cref{sec:sheafbw} are
formalized in Lean 4 (mathlib \texttt{v4.30.0-rc2}) and machine-verified
by the Lean kernel: \texttt{lake build} passes the entire library with
\emph{0 sorry, 0 admitted, 0 warnings} (8326 build jobs). Each formalized
result carries a footnote identifying the corresponding Lean file and
theorem name(s). The category-theoretic subsection
\cref{subsec:sheafbw-unitary-subcat} (i.e.,
\cref{def:unitary-subcat,prop:unitary-groupoid}), the physical instances
of \cref{sec:three_instances} (i.e.,
\cref{def:sampling-map,def:integration-map} and \cref{empirical:8b}), and
the geometric ``scalar multiple of isometry'' clause of
\cref{thm:strict-equality} are \emph{not} formalized; each such statement
carries an explicit ``Not formalized'' footnote at its declaration. The
complete paper--Lean correspondence table is in \texttt{lean/CROSSREF.md}.


% =============================================================
% Section 2. Related Work
% =============================================================
\section{Related Work}\label{sec:related}
% =============================================================
% Section 2. Related Work
% Template: paper/v0.1/notes/narrative.md §4 §2
% Structure: §2.1 FEEC / §2.2 cellular sheaf SP / §2.3 TopSP / §2.4 dagger Galerkin
% Key message: each school is precise; the four have not been combined.
% =============================================================

\noindent
The technical work in this paper combines tools from four schools
that have developed independently. Each school is internally precise;
none has been combined with the others on a single concrete object
producing a single scalar. We review the relevant results here so
that later sections can cite them by tag.

% -------------------------------------------------------------
\subsection{Finite element exterior calculus}\label{subsec:related-feec}

Finite element exterior calculus rests on three lines of
construction. The reduction--reconstruction interface of
\cite{BochevHyman2006} (with antecedents in
\cite{Whitney1957,Bossavit1988}) defines a map
$R : \Lambda^k(\Omega) \to C^k$ that integrates a smooth $k$-form
over chains, and an inverse $I_W : C^k \to \Lambda^k$ that lifts
chains to forms via Whitney shape functions; these two maps satisfy a
commuting Stokes identity (\cref{prop:commuting-stokes}). The
finite-element-system framework of \cite{Christiansen2010} packages
these data abstractly: an FES is a contravariant functor on a
category of cellular complexes, and Whitney forms are recovered as
its canonical instance. The Hilbert complex framework of
\cite{AFW2010,FalkWinther2014} (refining
\cite{BruningLesch1992,AFW2006}) supplies the analytic core that we
need: closed cochain complexes of Hilbert spaces, Hodge decomposition
with stalk-wise harmonic representatives, and bounded cochain
projections $\pi_h$ commuting with the differential up to a fixed
constant. The recent survey of \cite{ChristiansenRapetti2025}
collates 30 years of FEEC development and is one of our two anchors
for placing FEEC tools relative to the topological-signal-processing
literature.

% -------------------------------------------------------------
\subsection{Cellular sheaves and sheaf signal processing}\label{subsec:related-sheafsp}

Cellular sheaves of vector spaces, with cohomology computed by a
chain complex on a regular cell complex, are surveyed in
\cite[\S 6]{Robinson2014}. Robinson defines sheaf bandwidth as a
real-valued analytic invariant and closes the chapter with the open
question on its categorically intrinsic formulation that motivates
the present paper. An earlier sheaf-theoretic perspective on
sampling, complementary in topic to the bandwidth question studied
here, appears in \cite{Robinson2014Sampling}; the abstract sampling
cell complex $X^A$ and the strict-equality construction of
\cref{thm:strict-equality} are independent of that work, supplying a
specific finite-dimensional sampling sheaf on which Robinson's open
question admits a concrete spectral resolution. The spectral theory of cellular sheaf Laplacians
$L^k = (\delta^k)^* \delta^k + \delta^{k-1} (\delta^{k-1})^*$ is
developed in \cite{HansenGhrist2019}; we use Definition~1 of that
paper, especially Cond.~4 (boundaries of $1$-cells consist of two
distinct $0$-cells), to certify regularity of the abstract sampling
cell complex of \cref{def:Xsamp}. The conference
version \cite{HansenGhrist2019Allerton} introduces the spectral gap
quantities $\gamma_+, \Gamma_+$ on which our intrinsic invariant
$\sheafBW := \sqrt{\gamma_+/\Gamma_+}$ is built.

% -------------------------------------------------------------
\subsection{Topological signal processing}\label{subsec:related-topsp}

Topological signal processing on simplicial complexes is the program
initiated by \cite{BarbarossaSardellitti2020}: signals live on
$k$-cochains, the simplicial Laplacian replaces the graph Laplacian,
and bandlimited subspaces are defined as eigenspaces of the
simplicial Laplacian below a fixed threshold $B$. The school treats
sampling and reconstruction via Riesz frame inequalities in the
bandlimited subspace, an angle that we recover as condition~(b) of
\cref{thm:triangle-equiv}. The recent survey of
\cite{CerejeirasKahlerLegatiuk2025} is our second anchor for
placing topological signal processing relative to FEEC; that survey
does not invoke AFW Hilbert complexes.

% -------------------------------------------------------------
\subsection{Dagger Galerkin theory}\label{subsec:related-dagger}

A categorical formulation of Galerkin theory in which monomorphisms
carry a $\dagger$ (right-adjoint) structure is given by
\cite{LahtinenStenvall2020}; the abstract dagger setting is
\cite{HeunenVicary2019}. We use this thread sparingly: the dagger
tools clarify which subcategories of $\CellSh^{\fin}_{\Hilb}(X)$
carry a unitary structure, and they make
\cref{op:non-trivial-dagger} a precise question about the existence
of a non-trivial $\dagger$ on the full category. The unitary
subcategory $\CellSh^{\fin}_{\Hilb}(X)^{\mathrm{u}}$ of
\cref{def:unitary-subcat} is a groupoid
(\cref{prop:unitary-groupoid}); the question is whether this
restricted structure extends.

% -------------------------------------------------------------
\subsection{What has not been combined}\label{subsec:related-gap}

Each of the four schools above provides precise tools, and each
admits the underlying combinatorial object---a regular cell complex
with cochain spaces and a coboundary operator---without difficulty.
Yet the four have not been combined on a single concrete object
producing a single scalar that is at once intrinsic to the abstract
structure and computable from acquisition data. The references
cited in this section are mutually disjoint: \cite{AFW2010} does
not invoke cellular sheaves; \cite{HansenGhrist2019} does not
invoke FEEC; \cite{BarbarossaSardellitti2020} does not invoke
either; \cite{LahtinenStenvall2020} does not invoke any signal
processing applications. This paper supplies the missing object
(the abstract sampling cell complex $X^A$ of \cref{def:Xsamp}) and
the missing scalar (the strict equality of
\cref{thm:strict-equality}). \Cref{tab:feec-four-schools} in
\cref{sec:feec_review} records the four-school comparison once the
needed FEEC machinery has been reviewed.

\medskip\noindent
Each school provides precise tools but operates in isolation. We now
exhibit the first formal bridge between three of them via a triangle
equivalence theorem (\cref{sec:triangle}), which pinpoints exactly
where Robinson's open problem lives.


% =============================================================
% Section 3. Triangle Equivalence
% =============================================================
\section{Triangle Equivalence: FEEC Geometry, Hodge Spectrum, and Robinson Sheaf Cohomology}\label{sec:triangle}
% =============================================================
% Section 3. Triangle Equivalence: FEEC Geometry, Hodge Spectrum,
% and Robinson Sheaf Cohomology
% =============================================================

% --- Opening hook (continuing from Section 2) ---
We now construct the bridge promised in \cref{sec:intro}: three
formulations of well-conditioning for $k$-space sampling, drawn from
three schools that have not previously been combined---finite element
exterior calculus (FEEC) singular value bounds, topological signal
processing (TopSP) Riesz frame conditions, and Robinson sheaf
cohomology with sheaf bandwidth---are formally equivalent under the
common finite-dimensional sampling-operator normal form introduced
in \cref{subsec:tri-setup}. The equivalence is asserted in this
finite-dimensional normal form; the original infinite-dimensional
or qualitative formulations of the three source literatures are not
themselves claimed to be equivalent outside this normalization (see
\cref{subsec:tri-edge}, (L1)). The equivalence is stated as
\cref{thm:triangle-equiv} and proved in
\cref{subsec:tri-ab,subsec:tri-ac}. The proof of $(a) \Leftrightarrow
(c)$ produces a candidate definition of sheaf bandwidth that is not
categorically intrinsic; the categorically intrinsic version is
constructed in \cref{sec:sheafbw}. We close the section by formalizing
Robinson's open problem from \cite[\S 6]{Robinson2014} as a precisely
narrowed two-clause question (\cref{subsec:tri-narrowing}).

% =============================================================
\subsection{Setup and notation}\label{subsec:tri-setup}
% =============================================================

The well-conditioning of $k$-space sampling admits three equivalent
readings: an SVD-style ratio $\sigma_{\min}(A)/\sigma_{\max}(A)$ of the
sampling operator; a Riesz frame inequality on a bandlimited subspace;
and a sheaf-cohomological condition on an abstract sampling cell complex.
\Cref{thm:triangle-equiv} makes this precise: the three readings agree
on the same threshold $\tau \in (0,1]$. The setup below introduces the
common notation in which all three formulations live---the bandlimited
subspace $V_B$, the sampling operator $A$, and the condition number
$\kappa(A)$---before the triangle statement in
\cref{subsec:tri-theorem}.

Let $\Omega_k := B(0, k_{\max}) \subset \mathbb{R}^d$ be a closed ball
in $k$-space, and let $V = \{k_1, \ldots, k_N\} \subset \Omega_k$ be a
finite set of $N$ candidate sample positions. We do not impose
geometric structure on $V$ at this stage; the dependence of the
constructions on $V$ is recorded explicitly via $V_B(V)$ below and
discussed in \cref{sec:three_instances}.

Let $X(V)$ denote the Delaunay simplicial complex on $V$ (cells:
vertices $V$, edges $E$, faces $F$, tetrahedra $T$), oriented in the
sense of \cite[\S 2.1]{Hatcher2002}. Edge weights are
$w_{ij} := \lvert k_i - k_j \rvert^{-2}$ (with the convention
$w_{ij} = 0$ if $\{i,j\} \notin E$), and
$W := (w_{ij})_{1 \le i, j \le N}$. Set
$D := \mathrm{diag}(d_1, \ldots, d_N)$ with $d_i := \sum_j w_{ij}$,
and define the graph Laplacian
$L_0 := D - W \in \mathbb{R}^{N \times N}$. The eigenpairs of $L_0$
are denoted $(\mu_l, v_l)_{l = 0, \ldots, N-1}$ with
$0 = \mu_0 \le \mu_1 \le \cdots \le \mu_{N-1}$ and
$\{v_l\}_{l=0}^{N-1}$ an orthonormal basis of $\C^N$.

For a fixed bandwidth threshold $B \in (0, \mu_{N-1})$, define the
$B$-bandlimited subspace
\begin{equation}\label{eq:VB-def}
  V_B(V) := \mathrm{span}\{\, v_l : \mu_l \le B \,\} \;\subset\; \C^N,
  \qquad K := \dim V_B(V).
\end{equation}
We write $V_B$ for $V_B(V)$ when $V$ is fixed and emphasize the
$V$-dependence only when comparing across instances
(\cref{sec:three_instances}).

Fix a sampling subset $V_s \subset \{1, \ldots, N\}$ with
$M := \lvert V_s \rvert \ge K$. Let
$\Sigma_{V_s} : \C^N \to \C^M$ be the index restriction operator
$(\Sigma_{V_s} c)_n := c_{i_n}$ for $V_s = \{i_1, \ldots, i_M\}$, and
let $\iota_B : V_B \hookrightarrow \C^N$ be the canonical isometric
embedding. The \emph{sampling operator}
\begin{equation}\label{eq:A-def}
  A \;:=\; \Sigma_{V_s} \circ \iota_B \;:\; V_B \longrightarrow \C^M
\end{equation}
is the central object of this paper. In the basis
$\{v_l\}_{l=0}^{K-1}$ of $V_B$ and the canonical basis of $\C^M$, $A$
has the matrix representation
\[
  A \;=\; \bigl[\, v_l(i_n) \,\bigr]_{n = 1, \ldots, M;\, l = 0, \ldots, K-1}
  \;\in\; \C^{M \times K}.
\]

The singular values of $A$ are denoted
\[
  \sigma_{\max}(A) \;:=\; \sigma_0(A)
  \;\ge\; \sigma_1(A) \;\ge\; \cdots \;\ge\; \sigma_{K-1}(A)
  \;=:\; \sigma_{\min}(A) \;\ge\; 0.
\]
The condition number is
$\kappa(A) := \sigma_{\max}(A) / \sigma_{\min}(A) \in [1, \infty]$,
with $\kappa(A) = \infty$ when $\sigma_{\min}(A) = 0$ (rank
deficiency). Throughout, $\tau \in (0, 1]$ denotes a fixed
well-conditioning threshold; $A$ is \emph{$\tau$-well-conditioned}
when $\kappa(A) \le 1/\tau$, equivalently
$\sigma_{\min}(A) \ge \tau \, \sigma_{\max}(A)$.

% =============================================================
\subsection{Triangle equivalence theorem}\label{subsec:tri-theorem}
% =============================================================

\begin{theorem}[Triangle Equivalence]\label{thm:triangle-equiv}\footnote{Formalized in Lean 4 as \texttt{triangle\_equivalence} in \texttt{lean/MRI/Cellular/TriangleEquivalence.lean}. The Lean statement formalizes the finite-dimensional operator/sheaf core of the theorem, with $V_B$ represented by an abstract Hilbert space and $A$ by an evaluation operator; the construction of $V_B(V)$ from the Delaunay graph is treated as setup data rather than separately formalized geometry. Full correspondence: \texttt{lean/CROSSREF.md}.}
Fix $V$, $V_s$, $B$ (and hence $V_B = V_B(V)$ and $A$ as in
\cref{subsec:tri-setup}), and fix $\tau \in (0, 1]$. The following
three statements are equivalent.
\begin{enumerate}
  \item[\textup{(a)}] \emph{(FEEC singular value form.)}
        $\sigma_{\min}(A) \ge \tau \, \sigma_{\max}(A)$.
  \item[\textup{(b)}] \emph{(Hodge-spectrum Riesz form.)} There exist
        constants $0 < \alpha \le \beta$ with $\beta / \alpha \le 1/\tau$
        such that for every $c \in V_B \setminus \{0\}$,
        \begin{equation}\label{eq:riesz}
          \alpha \,\lVert c \rVert_{\C^N}
          \;\le\;
          \lVert \Sigma_{V_s} c \rVert_{\C^M}
          \;\le\;
          \beta \,\lVert c \rVert_{\C^N}.
        \end{equation}
        That is, $V_B$ forms a Riesz sequence in $\ell^2(V_s)$ with
        frame bound ratio $\beta / \alpha \le 1/\tau$.
  \item[\textup{(c)}] \emph{(Robinson sheaf form.)} On the
        \emph{abstract sampling sheaf} $\samps$ over the abstract
        sampling cell complex $X^A$, constructed in
        \cref{def:Xsamp,def:sampling-sheaf} of \cref{sec:sheafbw},
        \begin{equation}\label{eq:robinson-sheaf-form}
          H^0\!\left( X^A;\, \samps \right) \;=\; 0
          \qquad \text{and} \qquad
          \sheafBW^{\mathrm{cand}}\!\left( \samps \right)
          \;\ge\; \tau,
        \end{equation}
        where $\sheafBW^{\mathrm{cand}}$ is the candidate sheaf
        bandwidth of \cref{def:sheafbw-candidate} below.
\end{enumerate}
\end{theorem}

\begin{remark}[Three readings of one inequality]\label{rem:three-readings}
The three forms in \cref{thm:triangle-equiv} read the same scalar
inequality in three vocabularies developed in disjoint literatures:
the FEEC singular value form is the standard reading via singular
value bounds \cite{AFW2010}; the Riesz form is the TopSP reading via
$\ell^2$ frame bounds \cite{BarbarossaSardellitti2020}; the Robinson
form translates the inequality into the language of cellular sheaves
on $X(V)$ in the sense of \cite{Robinson2014}, with the sheaf
bandwidth furnishing the quantitative bound.
\end{remark}

\Cref{thm:triangle-equiv} is proved in two parts:
$(a) \Leftrightarrow (b)$ in \cref{subsec:tri-ab}, and
$(a) \Leftrightarrow (c)$ in \cref{subsec:tri-ac}. The latter relies
on the abstract sampling sheaf $\samps$ over the cell complex $X^A$,
whose explicit construction is given in
\cref{subsec:sheafBW-Xsamp} of \cref{sec:sheafbw}; the candidate
sheaf bandwidth $\sheafBW^{\mathrm{cand}}$ used in clause~\textup{(c)}
is recorded in \cref{def:sheafbw-candidate} below in terms of
$\samps$.

% =============================================================
\subsection{Proof of $(a) \Leftrightarrow (b)$}\label{subsec:tri-ab}
% =============================================================

The equivalence $(a) \Leftrightarrow (b)$ is a direct consequence of
the variational characterization of singular values together with the
isometry of $\iota_B$.

\begin{proof}[Proof of $(a) \Leftrightarrow (b)$]
Since $\iota_B : V_B \hookrightarrow \C^N$ is an isometric embedding,
$\lVert \iota_B c \rVert_{\C^N} = \lVert c \rVert_{V_B}$ for every
$c \in V_B$, and consequently
\begin{equation}\label{eq:Avariational}
  \lVert A c \rVert_{\C^M}
  \;=\;
  \lVert \Sigma_{V_s} \iota_B c \rVert_{\C^M}
  \;=\;
  \lVert \Sigma_{V_s} c \rVert_{\C^M}
  \qquad \text{for every } c \in V_B,
\end{equation}
where in the last equality $V_B$ is identified with its image
$\iota_B(V_B) \subset \C^N$. The singular values of $A$ are then
characterized variationally by
\begin{align}
  \sigma_{\max}(A)
    &\;=\; \max_{\substack{c \in V_B \\ \lVert c \rVert = 1}}
        \lVert \Sigma_{V_s} c \rVert_{\C^M},
    \label{eq:sigma-max-var} \\
  \sigma_{\min}(A)
    &\;=\; \min_{\substack{c \in V_B \\ \lVert c \rVert = 1}}
        \lVert \Sigma_{V_s} c \rVert_{\C^M}.
    \label{eq:sigma-min-var}
\end{align}

\paragraph{$(b) \Rightarrow (a)$.}
Given $\alpha, \beta$ with the bounds \cref{eq:riesz}, take infimum
and supremum over $c \in V_B$ with $\lVert c \rVert = 1$. From
\cref{eq:sigma-max-var,eq:sigma-min-var} this yields
$\sigma_{\max}(A) \le \beta$ and $\sigma_{\min}(A) \ge \alpha$. The
hypothesis $\beta / \alpha \le 1/\tau$ then gives
$\sigma_{\min}(A) / \sigma_{\max}(A) \ge \alpha / \beta \ge \tau$,
which is $(a)$.

\paragraph{$(a) \Rightarrow (b)$.}
Set $\alpha := \sigma_{\min}(A)$ and $\beta := \sigma_{\max}(A)$.
Then $\beta / \alpha = \kappa(A) \le 1/\tau$ by hypothesis.
\Cref{eq:sigma-max-var,eq:sigma-min-var} together with
\cref{eq:Avariational} yield, after homogeneity in $c$, the bound
\cref{eq:riesz} with these constants.
\end{proof}

% =============================================================
\subsection{Proof of $(a) \Leftrightarrow (c)$}\label{subsec:tri-ac}
% =============================================================

The equivalence $(a) \Leftrightarrow (c)$ is realized through the
abstract sampling sheaf $\samps$ on the cell complex $X^A$, whose
explicit construction we forward-reference to
\cref{def:Xsamp,def:sampling-sheaf} in
\cref{subsec:sheafBW-Xsamp}. Two facts about $\samps$ proved there
suffice for the present argument:
\begin{enumerate}[label=\textup{(F\arabic*)},leftmargin=*]
  \item \emph{Cochain spaces.} $C^0(\samps) \cong V_B$ (canonical
        isometry, \cref{lem:canonical-isometry}) and
        $C^1(\samps) = \C^M$
        (\cref{lem:samps-cochain-spaces}).
  \item \emph{Coboundary equals sampling operator.} The $0$-th
        coboundary $\delta^0 : C^0(\samps) \to C^1(\samps)$ of
        $\samps$ coincides with $A$ after the canonical
        identification of \cref{lem:canonical-isometry,lem:samps-cochain-spaces},
        i.e.\ $\delta^0 = A$ as linear maps $V_B \to \C^M$
        (\cref{prop:delta-equals-A}).
\end{enumerate}
Both facts depend only on the data of \cref{def:Xsamp,def:sampling-sheaf}
and on the sampling operator $A$ recorded in \cref{subsec:tri-setup};
neither depends on \cref{thm:triangle-equiv}, so the forward
references are not circular.

\begin{definition}\label{def:sheafbw-candidate}
The \emph{candidate sheaf bandwidth} of $\samps$ is
\begin{equation}\label{eq:sheafbw-candidate}
  \sheafBW^{\mathrm{cand}}(\samps)
  \;:=\;
  \frac{\sigma_{\min}(A)}{\sigma_{\max}(A)}
  \;\in\;
  [0, 1],
\end{equation}
with the convention
$\sheafBW^{\mathrm{cand}}(\samps) := 0$ when
$\sigma_{\min}(A) = 0$. This candidate references $\samps$ through
the operator $A$ derived from the chosen embedding
$V_B \hookrightarrow \C^N$ and the basis $\{v_l\}_{l=0}^{K-1}$; it is
therefore not categorically intrinsic to the cellular sheaf data of
$\samps$. The categorically intrinsic version is the spectral-gap
ratio
$\sheafBW := \sqrt{\gamma_+(\delta^0) / \Gamma_+(\delta^0)}$
constructed in \cref{sec:sheafbw}, which on $\samps$ satisfies a
strict identity with $\sigma_{\min}(A) / \sigma_{\max}(A)$ via
\cref{thm:strict-equality}.
\end{definition}

\begin{proof}[Proof of $(a) \Leftrightarrow (c)$]
By \textup{(F1)}, the coboundary $\delta^0$ of $\samps$ is well-typed
as a linear map $V_B \to \C^M$, so the identification with $A$ in
\textup{(F2)} is type-compatible. Combining \textup{(F1)} and
\textup{(F2)} gives $\delta^0 = A$ as maps $V_B \to \C^M$, hence
$H^0(\samps) = \ker \delta^0 = \ker A$, and so
$H^0(\samps) = 0$ if and only if $\sigma_{\min}(A) > 0$. By
\cref{def:sheafbw-candidate},
$\sheafBW^{\mathrm{cand}}(\samps) \ge \tau$ if and only if
$\sigma_{\min}(A) \ge \tau \, \sigma_{\max}(A)$. The conjunction of
the two conditions in \cref{eq:robinson-sheaf-form} is therefore
equivalent to $\sigma_{\min}(A) \ge \tau \, \sigma_{\max}(A)$, i.e.\
$(a)$.
\end{proof}

% =============================================================
\subsection{Reading the candidate bandwidth}\label{subsec:tri-bw-candidate}
% =============================================================

The candidate definition \cref{eq:sheafbw-candidate} is a numerical
quantity attached to the sheaf $\samps$ via the operator $A$. Three
observations describe its character.

\paragraph{(i) Range.}
$\sheafBW^{\mathrm{cand}}(\samps) \in [0, 1]$, with the value $1$
attained when $A$ is a positive scalar multiple of an isometry on
$V_B$ (the maximally well-conditioned case;
cf.\ \cref{thm:strict-equality}) and the value $0$ attained when $A$
is rank-deficient (failure of $H^0$ vanishing).

\paragraph{(ii) Three readings, one inequality.}
\Cref{rem:three-readings} identifies the three vocabularies in which
$\sigma_{\min}(A) / \sigma_{\max}(A) \ge \tau$ is read; the candidate
definition \cref{eq:sheafbw-candidate} translates the FEEC reading
verbatim into the sheaf-cohomological reading. The TopSP reading via
Riesz frame bounds is recovered through clause $(b)$ of
\cref{thm:triangle-equiv}.

\paragraph{(iii) Lack of categorical intrinsicness.}
\Cref{def:sheafbw-candidate} uses the basis $\{v_l\}_{l=0}^{K-1}$ of
$V_B$ implicit in the matrix form of $A$, i.e.\ data of the
representation $V_B \hookrightarrow \C^N$ rather than data intrinsic
to the cellular sheaf $\samps$. The categorically intrinsic
alternative, defined in \cref{sec:sheafbw} as a spectral gap ratio
of the sheaf coboundary operator, refers only to $\mathscr{A}$, not
to an ambient embedding or to a chosen basis of stalks.

% =============================================================
\subsection{Robinson's open problem: precise narrowing}\label{subsec:tri-narrowing}
% =============================================================

Robinson \cite[\S 6]{Robinson2014} writes, in the context of cellular
sheaves whose stalks are Hilbert spaces and whose cochain complexes
admit Hodge decompositions, that ``there may be a way to discuss the
concept of \emph{bandwidth} for general sheaves of Hilbert spaces.''
\Cref{thm:triangle-equiv} together with \cref{def:sheafbw-candidate}
narrows this to a two-clause open problem.

\begin{enumerate}
  \item[\textup{(R1)}] \emph{(Categorical intrinsicness.)} Define a
        sheaf bandwidth $\sheafBW(\mathscr{A})$ for cellular sheaves
        $\mathscr{A}$ of Hilbert spaces using only the cochain-complex
        data of $\mathscr{A}$ (stalks, restriction maps, coboundary),
        without reference to an ambient embedding
        $V_B \hookrightarrow \C^N$ or to a chosen basis of stalks.
  \item[\textup{(R2)}] \emph{(Reduction on the sampling instance.)}
        Whatever intrinsic definition is given in \textup{(R1)}, it
        must reduce to $\sigma_{\min}(A) / \sigma_{\max}(A)$ on the
        sampling sheaf instance, so that the candidate
        \cref{eq:sheafbw-candidate} is recovered as a special case.
\end{enumerate}

The narrowing replaces the qualitative ``find an equivalent
definition'' in \cite[\S 6]{Robinson2014} with a concrete two-clause
specification. \Cref{sec:sheafbw} answers \textup{(R1)} by the
spectral-gap ratio definition
$\sheafBW := \sqrt{\gamma_+(\delta^0) / \Gamma_+(\delta^0)}$, and
answers \textup{(R2)} by \cref{thm:strict-equality}, which establishes
the strict identity
$\sheafBW(\samps) = \sigma_{\min}(A) / \sigma_{\max}(A)$ on the
abstract sampling cell complex of \cref{sec:sheafbw}.

% =============================================================
\subsection{Edge of soundness}\label{subsec:tri-edge}
% =============================================================

\Cref{thm:triangle-equiv} is stated and proved within the
finite-dimensional setting of \cref{subsec:tri-setup}: $V$ is a
finite vertex set in $\Omega_k$, $V_B \subset \C^N$ is finite
dimensional, and $A$ acts between finite-dimensional Hilbert spaces.
Four boundary conditions delimit the scope of the theorem; each is
recorded as a project-internal limitation, fully resolved within
this paper or instantiated in a subsequent section.

\paragraph{(L1) Infinite-dimensional stalks.}
Robinson's open problem in \cite[\S 6]{Robinson2014} is stated for
sheaves of Hilbert spaces with potentially infinite-dimensional
stalks, in the setting of Hilbert complexes \cite{BruningLesch1992}.
The triangle equivalence is established here under the
finite-dimensional restriction $\dim V_B = K < \infty$. The
categorically intrinsic version of \cref{sec:sheafbw} is constructed
in the same finite-dimensional restriction, in which closedness of
cochain ranges is automatic
(see \cref{thm:closed-hilbert-complex}). The infinite-dimensional
extension is outside the scope of this paper.

\paragraph{(L2) Categorical intrinsicness of the bandwidth.}
\Cref{def:sheafbw-candidate} is not categorically intrinsic to the
sheaf $\samps$: the basis $\{v_l\}_{l=0}^{K-1}$ underlying the
matrix form of $A$ is not part of the cellular sheaf data. The
categorically intrinsic version is given in \cref{sec:sheafbw}; it
answers clause \textup{(R1)} of \cref{subsec:tri-narrowing} and is
shown to satisfy clause \textup{(R2)} on $\samps$ via
\cref{thm:strict-equality}.

\paragraph{(L3) Spectral gap as data, not hypothesis.}
The bandlimited subspace
$V_B(V) = \mathrm{span}\{v_l : \mu_l \le B\}$ presupposes a separating
gap in the spectrum of $L_0$ at $B$; otherwise the cutoff is not
stable under perturbations of $V$. In the finite-dimensional finite
vertex setting, the spectral gap is a property of the data $(V, B)$
rather than a hypothesis of \cref{thm:triangle-equiv}: the theorem
is stated for fixed $V, V_s, B$ and makes no genericity assumption
on the gap. The three instances of \cref{sec:three_instances}
instantiate the spectral-gap data concretely; the resulting
trade-off between sampling and integration is recorded in
\cref{sec:three_instances} as Empirical Proposition~8b.

\paragraph{(L4) Vertex geometry not made explicit.}
The vertex geometry of $V$ enters \cref{thm:triangle-equiv} only
through the dependence of $V_B(V)$ on $V$ (and through it the matrix
$A = [v_l(i_n)]$). The geometry is not an independent variable of
the theorem statement. The three instances of
\cref{sec:three_instances} instantiate the abstract setting in three
domain-specific geometries (simplicial complexes for TopSP,
finite-element meshes for FEM-EM, Brillouin-zone tetrahedra for
condensed matter), and Empirical Proposition~8b records the Pareto
trade-off between sampling and integration across geometries within
each instance.

% --- Closing hook (to Section 4) ---
\Cref{thm:triangle-equiv} establishes the bridge as a formal
equivalence and narrows Robinson's open problem to clauses
\textup{(R1)} and \textup{(R2)} of \cref{subsec:tri-narrowing}. The
remaining work---constructing the categorically intrinsic sheaf
bandwidth that answers \textup{(R1)}---requires the FEEC machinery,
which we review in \cref{sec:feec_review} before resuming the
construction in \cref{sec:sheafbw}.


% =============================================================
% Section 4. FEEC W^0 / P Review
% =============================================================
\section{Finite Element Exterior Calculus: $W^0$ Lift and $P$ Projection}\label{sec:feec_review}
% =============================================================
% Section 4. Finite Element Exterior Calculus: W^0 Lift and P Projection
%
% Character: comparative review chapter. We introduce no new theory.
% =============================================================

% --- Opening hook (continuing from Section 3) ---
The technical core of Robinson's open problem---a sheaf-categorically
intrinsic bandwidth---requires lifting Hilbert-complex structure,
Hodge decomposition, and bounded cochain projection from the
continuous setting of finite element exterior calculus (FEEC) to
cellular sheaves of Hilbert spaces. We review the necessary FEEC
machinery in this section. The exposition is comparative: four
schools---mimetic discretization \cite{BochevHyman2006}, finite
element systems \cite{Christiansen2010}, FEEC Hilbert complexes
\cite{AFW2010}, and dagger Galerkin \cite{LahtinenStenvall2020}---
have developed the same family of constructions in disjoint
vocabularies over three decades. \emph{This section is a
comparative review; we do not introduce new theory.} Its purpose is
disambiguation and a side-by-side mapping that makes the
sheaf-side interface of \cref{sec:sheafbw} statable.

% =============================================================
\subsection{Bochev--Hyman reduction and reconstruction}\label{subsec:feec-bh}
% =============================================================

The mimetic-discretization framework of Bochev--Hyman
\cite{BochevHyman2006} organizes the connection between continuous
differential forms and discrete cochains around a pair of maps and a
commuting diagram.

\begin{definition}[Reduction and Whitney reconstruction]\label{def:bh-RI}
Let $\mathcal{T}$ be a cellular complex and $\Lambda^k(\Omega)$ the
space of differential $k$-forms on a domain $\Omega \supset
\lvert \mathcal{T} \rvert$. The \emph{reduction} (or de Rham map)
\[
  R : \Lambda^k(\Omega) \longrightarrow C^k(\mathcal{T}),
  \qquad
  R(\omega)(\sigma) := \int_\sigma \omega
  \quad\text{for $k$-cells } \sigma \in \mathcal{T},
\]
is the integration cochain map \cite[\S 1]{BochevHyman2006}, and the
\emph{Whitney reconstruction}
\[
  I_W : C^k(\mathcal{T}) \longrightarrow V_W^k \subset \Lambda^k(\Omega),
\]
sends a cochain to the unique Whitney $k$-form interpolating it on
each cell of $\mathcal{T}$
\cite{Whitney1957,Bossavit1988}.\footnote{Bochev--Hyman~%
\cite{BochevHyman2006} originally write $R/I$ for the
reduction/reconstruction pair. To avoid clash with the integration
map $I$ used in \cref{sec:three_instances} for the
sampling/integration duality, we write $R/I_W$ throughout this paper,
with subscript $W$ indicating Whitney reconstruction.}
\end{definition}

\begin{proposition}[Commuting Stokes; cf.\ {\cite[\S 2]{BochevHyman2006}}]
\label{prop:commuting-stokes}
With $R$ and $I_W$ as above, $R \circ d = \delta \circ R$, where $d$
is the exterior derivative on $\Lambda^k(\Omega)$ and $\delta$ is the
cochain coboundary on $C^k(\mathcal{T})$. Equivalently, $R$ is a
chain map between the de Rham complex of $\Omega$ and the cochain
complex of $\mathcal{T}$. The natural inner product on
$C^k(\mathcal{T})$ induced by Whitney reconstruction is
$\langle c, c' \rangle_{C^k} := \langle I_W c, I_W c' \rangle_{L^2(\Omega)}$.
\end{proposition}

The pair $(R, I_W)$ behaves asymmetrically with respect to composition:

\begin{remark}[$I_W$ is not a projection]\label{rem:IW-not-projection}
On $V_W^k \subset \Lambda^k(\Omega)$ one has $R \circ I_W = \mathrm{id}_{C^k}$
\cite{Whitney1957}, but on the larger ambient $\Lambda^k(\Omega)$ the
composition $I_W \circ R$ is not the identity, and $I_W$ is not a
bounded $L^2$-projection onto $V_W^k$. The same observation is
recorded explicitly by Falk--Winther \cite{FalkWinther2014}: ``the
operator is not a projection since it is not equal to the identity on
the Whitney forms.'' The repair is the construction of a
\emph{bounded cochain projection} $\pi_h$ in
\cref{subsec:feec-afw}.
\end{remark}

% =============================================================
\subsection{Christiansen finite element systems}\label{subsec:feec-fes}
% =============================================================

The mimetic framework of \cref{subsec:feec-bh} acquires a contravariant
functor form in Christiansen's finite element system (FES) framework
\cite{Christiansen2010}.

\begin{definition}[Finite element system; {\cite[\S 2, Def.~2]{Christiansen2010}}]
\label{def:christiansen-fes}
Let $\mathcal{T}$ be a cellular complex. A \emph{finite element system}
on $\mathcal{T}$ assigns, to each $k \in \mathbb{N}$ and each cell
$T \in \mathcal{T}$, a vector space
$\mathcal{E}^k(T) \subseteq \widehat{X}^k(T)$
of \emph{differential $k$-elements}\footnote{We follow
\cite[\S 2, Def.~2]{Christiansen2010} with one notational change: we
write $\mathcal{E}^k(T)$ for what \cite{Christiansen2010} denotes
$A^k(T)$, since $A$ is reserved throughout this paper for the
sampling operator $A : V_B \to \C^M$ of \cref{eq:A-def}.}
such that
\begin{enumerate}
  \item $d : \mathcal{E}^k(T) \to \mathcal{E}^{k+1}(T)$ is well
        defined; and
  \item for every cell inclusion $i : T' \subseteq T$, the pullback
        $i^\star : \mathcal{E}^k(T) \to \mathcal{E}^k(T')$ is a
        morphism of vector spaces.
\end{enumerate}
\end{definition}

\Cref{def:christiansen-fes} exhibits an FES as a contravariant functor
from the poset category of $\mathcal{T}$ (objects: cells; morphisms:
inclusions) to $\Vect$, with $T \mapsto \mathcal{E}^k(T)$ and
inclusions sent to pullbacks. Christiansen and Rapetti
\cite{ChristiansenRapetti2025} state this categorical reading
explicitly: ``a FES is just a contravariant functor from the
(poset) category determined by the mesh, to the category of vector
spaces.'' Whitney $k$-forms are an instance:
\cite[Example~2.3]{Christiansen2010} identifies the case of the
trimmed polynomial complex of degree~$1$ with the Whitney form
construction of \cite{Whitney1957}, applied in computational
electromagnetics by \cite{Bossavit1988}.

The de Rham map of an FES recovers cohomology in the following sense.

\begin{proposition}[{\cite[Prop.~2.10]{Christiansen2010}}]\label{prop:fes-derham}
Under the compatibility hypotheses of
\cite[\S 2, Def.~3]{Christiansen2010}, the de Rham map from an FES to
the simplicial cochain complex of $\mathcal{T}$ induces an isomorphism
on cohomology.
\end{proposition}

% =============================================================
\subsection{AFW Hilbert complexes and bounded cochain projection}\label{subsec:feec-afw}
% =============================================================

The framework of \cref{subsec:feec-bh,subsec:feec-fes} acquires its
analytic completion in the FEEC Hilbert-complex framework of
Arnold--Falk--Winther \cite{AFW2006,AFW2010}, which extends the
algebraic structure to closed densely-defined operators between
Hilbert spaces.

\begin{definition}[Hilbert complex; {\cite[\S 3.1.3]{AFW2010},
\cite{BruningLesch1992}}]\label{def:afw-hilbert-complex}
A \emph{Hilbert complex} is a sequence $(W^k, d^k)_{k \ge 0}$ in
which each $W^k$ is a Hilbert space and each
$d^k : W^k \to W^{k+1}$ is a closed, densely defined linear operator
satisfying $\mathrm{range}(d^k) \subseteq \mathrm{dom}(d^{k+1})$ and
$d^{k+1} \circ d^k = 0$. The complex is \emph{closed} when, for each
$k$, the range of $d^k$ is closed in $W^{k+1}$.
\end{definition}

Closedness yields the standard Hodge decomposition. Writing
$(d^k)^*$ for the Hilbert adjoint of $d^k$, set
$\Delta_k := d^{k-1} (d^{k-1})^* + (d^k)^* d^k$.

\begin{theorem}[Hodge decomposition;
{\cite[\S 3.2]{AFW2010},\cite{BruningLesch1992}}]\label{thm:afw-hodge}
For a closed Hilbert complex $(W^k, d^k)$,
\[
  W^k \;=\; \overline{\mathrm{range}(d^{k-1})}
  \;\oplus\; \mathrm{ker}(\Delta_k)
  \;\oplus\; \overline{\mathrm{range}((d^k)^*)},
\]
the three summands being mutually orthogonal. The harmonic space
$\mathcal{H}^k := \mathrm{ker}(\Delta_k)$ is canonically isomorphic
to the cohomology $H^k$ of the complex.
\end{theorem}

The discrete-side structure that connects $W^k$ to a finite-dimensional
subcomplex $W_h^k$ is the \emph{bounded cochain projection}.

\begin{definition}[Bounded cochain projection;
{\cite[\S 3.1.1]{AFW2010},\cite{FalkWinther2014}}]\label{def:bounded-pi}
A family of operators $\pi_h^k : W^k \to W_h^k$ is a
\emph{bounded cochain projection} when, for each $k$, it satisfies
\begin{enumerate}
  \item $\pi_h^k \circ \iota_h = \mathrm{id}_{W_h^k}$, where
        $\iota_h : W_h^k \hookrightarrow W^k$ is the canonical
        inclusion (true projection);
  \item $d^k \circ \pi_h^k = \pi_h^{k+1} \circ d^k$ on
        $\mathrm{dom}(d^k)$ (commutation with the coboundary);
  \item $\lVert \pi_h^k \rVert$ is uniformly bounded with respect to
        the mesh parameter $h$.
\end{enumerate}
\end{definition}

\begin{theorem}[{\cite[Thm.~3.6]{AFW2010}}]\label{thm:afw-cohomology}
If $(W_h^k, d^k)$ is a subcomplex of a closed Hilbert complex
$(W^k, d^k)$ admitting a bounded cochain projection, then the
cohomology of the subcomplex is canonically isomorphic to that of
the ambient complex, and the harmonic spaces have equal dimension.
\end{theorem}

Falk--Winther \cite{FalkWinther2014} construct a \emph{local}
bounded cochain projection meeting all three conditions of
\cref{def:bounded-pi} on simplicial subcomplexes; the construction
repairs the obstruction in \cref{rem:IW-not-projection} by replacing
$I_W$ with a smoothed extension that retains the commutation with
$d$ while becoming a true projection. Christiansen
\cite[\S 4]{Christiansen2010} gives a related family of
quasi-interpolators.

% =============================================================
\subsection{Lahtinen--Stenvall dagger mono form of Galerkin}\label{subsec:feec-dagger}
% =============================================================

A complementary categorical reading is provided by Lahtinen and
Stenvall \cite{LahtinenStenvall2020}, who characterize Galerkin
discretization as a morphism in the dagger category
$\Hilb$ of Hilbert spaces with bounded linear maps.

A \emph{dagger category} \cite{HeunenVicary2019} is a category
equipped with an involutive contravariant identity-on-objects functor
$\dagger$; in $\Hilb$, $f^\dagger$ is the Hilbert adjoint of
$f : H_1 \to H_2$. A \emph{dagger mono} is a morphism
$\iota : H_h \to H$ with $\iota^\dagger \circ \iota = \mathrm{id}_{H_h}$,
equivalently an isometric embedding.

\begin{proposition}[Galerkin as dagger mono;
{\cite{LahtinenStenvall2020}}]\label{prop:galerkin-dagger}
Discretization of a finite element method, viewed as a morphism in
$\Hilb$, is a dagger mono $\iota_h : V_h \hookrightarrow V$ from a
finite-dimensional discrete domain into the continuous Hilbert
space. The Galerkin projection $V \to V_h$ is the dagger
$\iota_h^\dagger$ of this embedding.
\end{proposition}

\Cref{prop:galerkin-dagger} aligns the FEEC framework with the
dagger-categorical formulation: $\iota_h$ is a dagger mono in
$\Hilb$ in the same sense as the inclusion of the Whitney form
subspace $V_W^k \hookrightarrow \Lambda^k(\Omega)$, and the Galerkin
projection is its Hilbert adjoint. The bounded cochain projection
$\pi_h$ of \cref{def:bounded-pi} is not, in general, a dagger
projection (i.e., it need not equal $\iota_h^\dagger$), since the
local construction of \cite{FalkWinther2014} sacrifices
self-adjointness in exchange for locality and uniform boundedness.

% =============================================================
\subsection{Four-school comparison table}\label{subsec:feec-comparison}
% =============================================================

The four schools of \cref{subsec:feec-bh,subsec:feec-fes,subsec:feec-afw,subsec:feec-dagger}
develop the same family of constructions in disjoint vocabularies.
\Cref{tab:feec-four-schools} records the side-by-side mapping.

\begin{table}[h]
  \centering
  \small
  \begin{tabular}{@{}llll@{}}
    \toprule
    Construction & Bochev--Hyman \cite{BochevHyman2006} &
       Christiansen \cite{Christiansen2010} &
       AFW \cite{AFW2010} \\
    \midrule
    Continuous space &
       $\Lambda^k(\Omega)$ & $\widehat{X}^k(T)$ & $W^k$ \\
    Discrete space &
       $C^k(\mathcal{T})$ & $\mathcal{E}^k(T)$ & $W_h^k$ \\
    Coboundary & $\delta$ on $C^k$ & $d$ on $\mathcal{E}^k$ &
       $d^k$ on $W^k$ \\
    Lift to ambient &
       $I_W$ (Whitney) & FES inclusions \cite[Ex.~2.3]{Christiansen2010} &
       $\iota_h : W_h^k \hookrightarrow W^k$ \\
    Reduction to discrete &
       $R$ (de Rham map) & pullbacks $i^\star$ & $\pi_h^k$
       (bounded cochain projection) \\
    Commutation &
       $R \circ d = \delta \circ R$ & functoriality of FES &
       $d^k \circ \pi_h^k = \pi_h^{k+1} \circ d^k$ \\
    Cohomology agreement &
       (implicit via Whitney 1957) & \cite[Prop.~2.10]{Christiansen2010} &
       \cite[Thm.~3.6]{AFW2010} \\
    \bottomrule
  \end{tabular}
  \caption{The mimetic, FES, and AFW frameworks for the
    same family of constructions, mapped column by column. The
    Lahtinen--Stenvall framework \cite{LahtinenStenvall2020} reads
    the inclusion row as a dagger mono in
    $\Hilb$.}\label{tab:feec-four-schools}
\end{table}

\begin{remark}[Common scope]\label{rem:feec-scope}
Each of the four schools establishes its constructions for a single
Hilbert complex (continuous side) related to a single subcomplex
(discrete side). None of the four schools attaches the constructions
to a sheaf over a cell complex with non-trivial restriction maps
between stalks. The transition to the sheaf side is taken in
\cref{sec:sheafbw}.
\end{remark}

% =============================================================
\subsection{Robinson's open problem: technical interface}\label{subsec:feec-robinson-interface}
% =============================================================

Robinson \cite[\S 6]{Robinson2014} writes that ``If $F$ and $S$ are
sheaves of Hilbert spaces then their cochain complexes are Hilbert
complexes, and they have a Hodge decomposition\ldots\ There may be
a way to discuss the concept of \emph{bandwidth} for general sheaves
of Hilbert spaces.'' \Cref{thm:triangle-equiv,def:sheafbw-candidate}
narrowed this to the two-clause problem
\textup{(R1)}--\textup{(R2)} of \cref{subsec:tri-narrowing}.
\Cref{sec:sheafbw} answers it in three steps, each of which uses one
ingredient of the FEEC machinery reviewed above.

\begin{enumerate}
  \item[\textbf{Step 1:}] \emph{Hilbert complex structure on cellular
        sheaves of Hilbert spaces.} Each stalk
        $\mathscr{A}(\sigma)$ becomes a Hilbert space, restriction
        maps become bounded operators, and the cochain spaces
        $C^k(\mathscr{A})$ assemble into a closed Hilbert complex
        in the sense of \cref{def:afw-hilbert-complex}. In the
        finite-dimensional setting of \cref{sec:sheafbw}, closedness
        of cochain ranges is automatic
        (\cref{thm:closed-hilbert-complex}). This step lifts
        \cref{def:afw-hilbert-complex} from a single chain to a
        sheaf-valued chain.
  \item[\textbf{Step 2:}] \emph{Hodge decomposition stalk-wise with
        restriction commutation.} \Cref{thm:afw-hodge} applied to
        the sheaf-cochain Hilbert complex of Step~1 yields a
        decomposition of $C^k(\mathscr{A})$ into exact, harmonic,
        and co-exact summands; the restriction maps of $\mathscr{A}$
        commute with the decomposition by functoriality. The
        sheaf Laplacian of \cite{HansenGhrist2019} is the
        finite-dimensional specialization.
  \item[\textbf{Step 3:}] \emph{Categorically intrinsic sheaf
        bandwidth.} \Cref{def:bounded-pi}'s view of
        well-conditioning as the boundedness of $\pi_h$ is replaced
        on the sheaf side by a spectral-gap ratio of the sheaf
        coboundary $\delta^0$, namely
        $\sheafBW := \sqrt{\gamma_+(\delta^0) /
        \Gamma_+(\delta^0)}$, defined intrinsically on the cellular
        sheaf data of $\mathscr{A}$. This step answers
        clauses~\textup{(R1)} and \textup{(R2)} of
        \cref{subsec:tri-narrowing}; it is constructed in
        \cref{sec:sheafbw}.
\end{enumerate}

Steps~1 and~2 are direct transcriptions of \cite{AFW2010}
\cref{def:afw-hilbert-complex,thm:afw-hodge} into the sheaf setting
and contribute no new theory. Step~3 is the categorically intrinsic
bandwidth definition; it is the locus of the new content of this
paper, completed in \cref{sec:sheafbw}.

% --- Closing hook (to Section 5) ---
\Cref{def:afw-hilbert-complex,thm:afw-hodge,def:bounded-pi}, together
with the FES contravariant functor form
\cref{def:christiansen-fes,prop:fes-derham} and the dagger
formulation \cref{prop:galerkin-dagger}, supply the algebraic
machinery: closed Hilbert complexes, Hodge decompositions, and
bounded cochain projections, with the categorical structures that
make their sheaf-side transcription unambiguous. What is missing on
the sheaf side---a categorically intrinsic bandwidth answering
clauses~\textup{(R1)} and~\textup{(R2)} of
\cref{subsec:tri-narrowing}---is constructed in
\cref{sec:sheafbw}.


% =============================================================
% Section 5. Sheaf Bandwidth Categorically Intrinsic
% =============================================================
\section{Sheaf Bandwidth: A Categorically Intrinsic Definition}\label{sec:sheafbw}
% =============================================================
% Section 5. Sheaf Bandwidth: A Categorically Intrinsic Definition
%
% Core technical chapter.
% =============================================================

% --- Opening hook (continuing from Section 4) ---
Building on the FEEC machinery of \cref{sec:feec_review}, we now
construct the missing piece: a categorically intrinsic candidate for
sheaf bandwidth, prove its invariance under unitary equivalence, and
establish a strict identity reducing it to a singular value ratio on
a canonical abstract sampling cell complex. The construction
answers clauses~\textup{(R1)} and~\textup{(R2)} of
\cref{subsec:tri-narrowing}: the definition refers only to the
cellular-sheaf data of $\mathscr{A}$
(\cref{def:sheafBW,prop:sheafBW-unitary}), and on the sampling sheaf
$\samps$ it agrees with $\sigma_{\min}(A)/\sigma_{\max}(A)$
strictly (\cref{thm:strict-equality}). Throughout this section all
stalks are finite dimensional; the infinite-dimensional extension
is outside the scope of this paper, in line with limitation~(L1) of
\cref{subsec:tri-edge}.

% =============================================================
\subsection{Key design choice: $\sheafBW := \sqrt{\gamma_+ / \Gamma_+}$}\label{subsec:sheafbw-design}
% =============================================================

The candidate of \cref{def:sheafbw-candidate} expressed sheaf
bandwidth via the basis-dependent quantity
$\sigma_{\min}(A)/\sigma_{\max}(A)$. We introduce a categorically
intrinsic alternative, formulated as a normalized spectral-gap
ratio of the sheaf coboundary.

\begin{definition}[Sheaf bandwidth]\label{def:sheafBW}
Let $\mathscr{A}$ be a finite-dimensional cellular sheaf of Hilbert
spaces on a finite regular cell complex $X$
(\cref{def:finite-dim-sheaf} below), and let $\delta^0$ be the
$0$-th coboundary of its cellular cochain complex
(\cref{def:coboundary}). The \emph{sheaf bandwidth} of
$\mathscr{A}$ is
\begin{equation}\label{eq:sheafBW-def}
  \sheafBW(\mathscr{A}) \;:=\;
  \sqrt{\frac{\gamma_+(\delta^0)}{\Gamma_+(\delta^0)}}
  \;\in\; [0, 1],
\end{equation}
where $\gamma_+(\delta^0)$ and $\Gamma_+(\delta^0)$ are the
infimum of the positive part of the spectrum and the supremum of
the spectrum of $(\delta^0)^* \delta^0$, defined formally in
\cref{def:spectral-gap}. By convention,
$\sheafBW(\mathscr{A}) := 0$ when $\delta^0 \equiv 0$ (the
degenerate case in which $\mathscr{A}$ has no coboundary action;
the zero map yields zero conditioning).
\end{definition}

The choice of \cref{eq:sheafBW-def} is governed by three
considerations.

\paragraph{(a) Square root.}
$\gamma_+$ and $\Gamma_+$ are squared singular values of
$\delta^0$. On the sampling sheaf $\samps$ of
\cref{def:sampling-sheaf} below,
$\gamma_+(\delta^0) = \sigma_{\min}(A)^2$ and
$\Gamma_+(\delta^0) = \sigma_{\max}(A)^2$
(\cref{thm:strict-equality}); the square root produces a
length-scale ratio rather than an area ratio, aligning the bandwidth
with the dimensionless conditioning $\sigma_{\min}/\sigma_{\max}$
that frame theory takes as the standard frame-bound ratio.

\paragraph{(b) Ratio.}
The ratio is scale invariant: global rescaling
$\mathscr{A}' = c \cdot \mathscr{A}$ by a single scalar
$c \in \mathbb{R} \setminus \{0\}$ applied to every stalk leaves
$\gamma_+ / \Gamma_+$ unchanged, so
$\sheafBW(\mathscr{A}') = \sheafBW(\mathscr{A})$. This is the
categorical requirement that bandwidth distinguish cellular sheaves
up to an overall scalar.

\paragraph{(c) Numerator and denominator.}
The numerator $\gamma_+$ measures the smallest non-trivial action
of $\delta^0$, increasing with well-conditioned sheaves; the
denominator $\Gamma_+$ normalizes the result to $[0, 1]$. The
maximally well-conditioned case ($\sheafBW = 1$) corresponds to
$\delta^0$ proportional to a partial isometry on
$(\ker \delta^0)^\perp$. On the sampling sheaf $\samps$ under
the hypothesis $H^0(\samps) = 0$ of \cref{thm:strict-equality},
$(\ker \delta^0)^\perp = V_B$ and this reduces to $A$ being a
positive scalar multiple of an isometry on $V_B$
(cf.\ \cref{thm:strict-equality}).

\Cref{tab:sheafbw-comparison} positions \cref{def:sheafBW} alongside
the spectral-gap measure of Hansen and Ghrist
\cite{HansenGhrist2019} and the Galerkin singular value spectrum of
Lahtinen and Stenvall \cite{LahtinenStenvall2020}.

\begin{table}[h]
  \centering
  \small
  \begin{tabular}{@{}lll@{}}
    \toprule
    Quantity & Form & Range / dimension \\
    \midrule
    $\sheafBW(\mathscr{A})$ (this paper) &
       $\sqrt{\gamma_+(\delta^0)/\Gamma_+(\delta^0)}$ &
       dimensionless, $[0, 1]$ \\
    Sheaf spectral gap \cite{HansenGhrist2019} &
       $\inf \spec(L^k) \cap (0, \infty)$ &
       absolute, $(0, \infty)$ \\
    Galerkin singular spectrum \cite{LahtinenStenvall2020} &
       $\spec(\delta^* \delta)^{1/2}$ &
       set-valued, singular values \\
    Frame bound ratio (frame theory) &
       $A/B$ on $V_B$ &
       dimensionless, $[0, 1]$ \\
    \bottomrule
  \end{tabular}
  \caption{Comparison of \cref{def:sheafBW} with related
    quantities. The sheaf bandwidth coincides with the frame-bound
    ratio of frame theory and is the dimensionless normalization of
    the spectral gap of \cite{HansenGhrist2019}.}\label{tab:sheafbw-comparison}
\end{table}

The remainder of this section makes \cref{def:sheafBW} formally
precise, proves its unitary invariance, and reduces it to
$\sigma_{\min}(A)/\sigma_{\max}(A)$ on a canonical sampling sheaf.

% =============================================================
\subsection{Sheaf Laplacian and spectral gap}\label{subsec:sheafbw-spectral-gap}
% =============================================================

We make precise the spectral quantities entering \cref{def:sheafBW}.

\begin{definition}[Sheaf Laplacian; cf.\ {\cite[\S 3]{HansenGhrist2019}},
\cite{AFW2010}]\label{def:sheaf-laplacian}
For a finite-dimensional cellular sheaf $\mathscr{A}$ on a finite
regular cell complex $X$, with cochain complex
$(C^\bullet(\mathscr{A}), \delta^\bullet)$ of
\cref{def:cochain-space,def:coboundary} below, the \emph{sheaf
Laplacian} of degree $k$ is
\[
  L^k(\mathscr{A}) \;:=\; \delta^{k-1} (\delta^{k-1})^*
  + (\delta^k)^* \delta^k
  \;:\;
  C^k(\mathscr{A}) \longrightarrow C^k(\mathscr{A}),
\]
with $\delta^{-1} := 0$. By \cref{thm:closed-hilbert-complex},
$L^k(\mathscr{A})$ is a finite-dimensional self-adjoint positive
operator.
\end{definition}

\begin{definition}[Spectral gap]\label{def:spectral-gap}
For $\mathscr{A}$ as in \cref{def:sheaf-laplacian} and any $k$ for
which $\delta^k \ne 0$,
\begin{align*}
  \gamma_+(\delta^k) &\;:=\;
    \inf\{\, \lambda \in \spec((\delta^k)^* \delta^k) :
              \lambda > 0 \,\}, \\
  \Gamma_+(\delta^k) &\;:=\;
    \sup \spec((\delta^k)^* \delta^k)
    \;=\; \lVert \delta^k \rVert_{\mathrm{op}}^2.
\end{align*}
When $\delta^k = 0$ the positive part of
$\spec((\delta^k)^* \delta^k)$ is empty and $\gamma_+(\delta^k)$ is
left undefined; the degenerate case is handled directly by the
convention in \cref{def:sheafBW}.
\end{definition}

\begin{lemma}[Well-definedness]\label{lem:gap-well-defined}
The quantities $\gamma_+(\delta^k)$ and $\Gamma_+(\delta^k)$ depend
only on the categorical data of $\mathscr{A}$ (stalks and
restriction maps); they are independent of any choice of basis on
the stalks.
\end{lemma}

\begin{proof}
The operator $(\delta^k)^* \delta^k$ is self-adjoint and positive on
the finite-dimensional Hilbert space $C^k(\mathscr{A})$. Since
$C^k(\mathscr{A})$ is finite-dimensional, the spectrum coincides with
the eigenvalue set of the matrix representation in any orthonormal
basis \cite[\S VI.1, \S VII.1]{ReedSimonI}, and this eigenvalue set is
invariant under unitary change of basis. Since the inner product on
$C^k(\mathscr{A})$ is determined stalk-wise by the Hilbert structure
on each $\mathscr{A}(\sigma)$ (and not by a chosen basis), the
spectrum---and hence $\gamma_+, \Gamma_+$---is determined by the
sheaf data alone.
\end{proof}

% =============================================================
\subsection{Unitary invariance}\label{subsec:sheafbw-unitary}
% =============================================================

The categorical-intrinsicness clause of \cref{lem:gap-well-defined}
extends to invariance under unitary equivalence in the category of
cellular sheaves of Hilbert spaces.

\begin{proposition}[Unitary invariance of $\sheafBW$]\label{prop:sheafBW-unitary}\footnote{Formalized in Lean 4 as \texttt{sheafBW\_eq} in \texttt{lean/MRI/Cellular/Unitary.lean}. The Lean statement proves the stalk-wise unitary-morphism core on the project formalization's \texttt{Complex1D} cellular-sheaf structure for $\delta^0$; the broader categorical packaging of finite regular cell complexes is kept informal in \cref{subsec:sheafbw-unitary-subcat}. Full correspondence: \texttt{lean/CROSSREF.md}.}
Let $\mathscr{A}, \mathscr{A}'$ be finite-dimensional cellular
sheaves of Hilbert spaces on $X$, and suppose there exists a
\emph{unitary} sheaf morphism $\Phi : \mathscr{A} \to \mathscr{A}'$,
i.e., a family of stalk-wise unitaries
$\Phi_\sigma : \mathscr{A}(\sigma) \to \mathscr{A}'(\sigma)$
satisfying the cellular-sheaf compatibility
$\Phi_\tau \circ \mathscr{A}(\sigma \le \tau) =
\mathscr{A}'(\sigma \le \tau) \circ \Phi_\sigma$ for every $\sigma
\le \tau$. Then
\begin{equation}\label{eq:sheafBW-unitary}
  \sheafBW(\mathscr{A}) \;=\; \sheafBW(\mathscr{A}').
\end{equation}
\end{proposition}

\begin{proof}
The family $\{\Phi_\sigma\}$ assembles into a unitary operator
$U : C^\bullet(\mathscr{A}) \to C^\bullet(\mathscr{A}')$ (a
direct sum of stalk-wise unitaries is unitary on the Hilbert direct
sum). The compatibility relation
$\Phi_\tau \circ \mathscr{A}(\sigma \le \tau) =
\mathscr{A}'(\sigma \le \tau) \circ \Phi_\sigma$ implies
$U \circ \delta^k_{\mathscr{A}} = \delta^k_{\mathscr{A}'} \circ U$ on
each $C^k$. Taking adjoints,
$U \circ (\delta^k_{\mathscr{A}})^* (\delta^k_{\mathscr{A}}) \circ U^* =
(\delta^k_{\mathscr{A}'})^* (\delta^k_{\mathscr{A}'})$ on
$C^k(\mathscr{A}')$. Unitary conjugation preserves the spectrum
\cite[\S VI.1]{ReedSimonI}, so
$\gamma_+(\delta^k_{\mathscr{A}}) =
\gamma_+(\delta^k_{\mathscr{A}'})$ and
$\Gamma_+(\delta^k_{\mathscr{A}}) =
\Gamma_+(\delta^k_{\mathscr{A}'})$ for every $k$, whence
\cref{eq:sheafBW-unitary}.
\end{proof}

\begin{remark}[Extent of invariance]\label{rem:sheafbw-extent}
\Cref{prop:sheafBW-unitary} establishes invariance under the
\emph{unitary subcategory} of cellular sheaves of Hilbert spaces
(\cref{def:unitary-subcat}). Whether the invariance extends to a
non-trivial dagger structure on the full category is recorded as an
independent open problem (\cref{op:non-trivial-dagger}).
\end{remark}

% =============================================================
\subsection{Closed Hilbert complex structure}\label{subsec:sheafbw-hilbert-complex}
% =============================================================

We now record the structural ingredients required by
\cref{def:sheaf-laplacian,def:spectral-gap}: the cellular cochain
complex of a finite-dimensional cellular sheaf is a closed Hilbert
complex in the sense of \cref{def:afw-hilbert-complex}.

\begin{definition}[Finite-dimensional cellular sheaf of Hilbert
  spaces; cf.\ {\cite[\S 2, Def.~4]{HansenGhrist2019}},
  \cite{Robinson2014}]\label{def:finite-dim-sheaf}
Let $X$ be a finite regular cell complex with face poset
$(P_X, \le)$. A \emph{finite-dimensional cellular sheaf of Hilbert
spaces} $\mathscr{A}$ on $X$ assigns
\begin{enumerate}[label=(\alph*)]
  \item to each cell $\sigma \in P_X$, a finite-dimensional Hilbert
        space $\mathscr{A}(\sigma)$;
  \item to each $\sigma \le \tau$ in $P_X$, a linear operator
        $\mathscr{A}(\sigma \le \tau) :
        \mathscr{A}(\sigma) \to \mathscr{A}(\tau)$ (automatically
        bounded, by finite-dimensionality \cite[\S III.5]{ReedSimonI});
\end{enumerate}
satisfying the functoriality $\mathscr{A}(\sigma \le \rho) =
\mathscr{A}(\tau \le \rho) \circ \mathscr{A}(\sigma \le \tau)$ for
$\sigma \le \tau \le \rho$.
\end{definition}

\begin{definition}[Cochain spaces and coboundary]\label{def:cochain-space}
For $\mathscr{A}$ as in \cref{def:finite-dim-sheaf}, the
\emph{cochain spaces} are
\[
  C^k(\mathscr{A}) \;:=\;
  \bigoplus_{\substack{\sigma \in P_X \\ \dim \sigma = k}}
  \mathscr{A}(\sigma),
\]
the Hilbert direct sums of the stalks of dimension $k$.
\end{definition}

\begin{definition}[Coboundary]\label{def:coboundary}
For $\mathscr{A}$ as in \cref{def:finite-dim-sheaf}, the
\emph{coboundary} $\delta^k : C^k(\mathscr{A}) \to C^{k+1}(\mathscr{A})$
is defined by
\[
  (\delta^k c)(\tau) \;:=\;
  \sum_{\substack{\sigma : \dim \sigma = k \\ \sigma < \tau}}
  [\sigma : \tau] \cdot \mathscr{A}(\sigma \le \tau)\!\left(c(\sigma)\right),
\]
for $c \in C^k(\mathscr{A})$ and $\tau$ a $(k+1)$-cell, where
$[\sigma : \tau] \in \{\pm 1\}$ are the incidence coefficients of
the oriented cell complex (see \cite[\S 2.1]{Hatcher2002}).
\end{definition}

\begin{theorem}[Closed Hilbert complex; cf.\ {\cite[\S 3.1.3]{AFW2010}}]
\label{thm:closed-hilbert-complex}\footnote{Formalized in Lean 4 at the \texttt{Complex1D} / $\delta^0$ layer: \texttt{coboundaryHilbert\_range\_isClosed} in \texttt{lean/MRI/Cellular/HilbertComplex.lean} proves condition (d), range closedness. Conditions (a)--(b) are finite-dimensional auto-derivations and condition (c) is vacuous in the one-dimensional complex used by the formalization; the all-$k$ finite-regular-cell-complex statement is not packaged as one Lean theorem. Full correspondence: \texttt{lean/CROSSREF.md}.}
For $\mathscr{A}$ a finite-dimensional cellular sheaf of Hilbert
spaces on a finite regular cell complex, the cochain complex
$(C^\bullet(\mathscr{A}), \delta^\bullet)$ is a closed Hilbert
complex in the sense of \cref{def:afw-hilbert-complex}: each
$C^k(\mathscr{A})$ is a finite-dimensional Hilbert space; each
$\delta^k$ is a bounded linear operator; $\delta^{k+1} \circ \delta^k = 0$;
and each $\mathrm{range}(\delta^k)$ is closed in $C^{k+1}(\mathscr{A})$.
\end{theorem}

\begin{proof}
Each $C^k(\mathscr{A})$ is a finite Hilbert direct sum of
finite-dimensional Hilbert spaces and is therefore a
finite-dimensional Hilbert space. Each $\delta^k$ is a linear
operator between finite-dimensional Hilbert spaces and is
automatically bounded \cite[\S III.5]{ReedSimonI}. The condition
$\delta^{k+1} \circ \delta^k = 0$ follows from the standard
incidence identity
$\sum_{\tau : \sigma < \tau < \rho} [\sigma : \tau][\tau : \rho] = 0$
\cite[\S 2.1]{Hatcher2002}. Range closedness is immediate: any
subspace of a finite-dimensional Hilbert space is closed.
\end{proof}

\begin{lemma}[Spectral form of $L^0$]\label{lem:L0-spectrum}\footnote{Formalized in Lean 4: the substantive operator claims correspond to three separate Lean theorems in \texttt{lean/MRI/Cellular/HilbertComplex.lean}: \texttt{sheafLaplacian0\_isSymmetric}, \texttt{sheafLaplacian0\_isPositive}, and \texttt{sheafLaplacian0\_ker\_eq}. The finite-spectrum and norm-bound consequences are standard finite-dimensional/self-adjoint/positive-operator facts supplied by mathlib rather than separate project theorems. Full correspondence: \texttt{lean/CROSSREF.md}.}
For $\mathscr{A}$ as in \cref{def:finite-dim-sheaf},
$L^0(\mathscr{A}) = (\delta^0)^* \delta^0$ is a self-adjoint
positive operator on $C^0(\mathscr{A})$. Its spectrum is finite,
$\spec(L^0(\mathscr{A})) \subset [0, \lVert \delta^0 \rVert_{\mathrm{op}}^2]$,
and $\ker L^0(\mathscr{A}) = \ker \delta^0 = H^0(\mathscr{A})$.
\end{lemma}

\begin{proof}
Self-adjointness and positivity follow from the form
$T^* T$ for any operator $T$. Finiteness of the spectrum follows
from finite-dimensionality. The kernel identity is the standard
fact that $\ker(T^* T) = \ker T$ on a Hilbert space.
\end{proof}

% =============================================================
\subsection{Abstract sampling cell complex and sampling sheaf}\label{subsec:sheafBW-Xsamp}
% =============================================================

The reduction of \cref{def:sheafBW} to
$\sigma_{\min}(A)/\sigma_{\max}(A)$ on a sampling instance requires
a cellular complex on which the sampling operator $A$ of
\cref{eq:A-def} arises as the coboundary $\delta^0$. We construct
this abstract complex and the associated sheaf, and verify that the
construction lands in the cellular-sheaf framework of
\cite{HansenGhrist2019,Robinson2014}.

\begin{definition}[Abstract sampling cell complex]\label{def:Xsamp}
Let $V_s = \{i_1, \ldots, i_M\}$. The \emph{abstract sampling cell
complex} $X^A$ is the finite regular cell complex with
\begin{itemize}[leftmargin=2em]
  \item two $0$-cells, $\sigma_B$ and $\sigma_\infty$;
  \item $M$ one-cells $\tau_1, \ldots, \tau_M$, with each $\tau_i$
        having boundary $\partial \tau_i = \{\sigma_B, \sigma_\infty\}$;
  \item incidence coefficients
        $[\sigma_B : \tau_i] = +1$, $[\sigma_\infty : \tau_i] = -1$
        (\cite[\S 2.1]{Hatcher2002}, with each $\tau_i$ oriented
        from $\sigma_B$ to $\sigma_\infty$).
\end{itemize}
The cell $\sigma_B$ is the \emph{principal $0$-cell}; the cell
$\sigma_\infty$ is an \emph{auxiliary basepoint} with no physical
content---its presence ensures that each $\tau_i$ has two distinct
endpoints, the regularity condition required by
\cref{rem:Xsamp-regularity} below.
\end{definition}

\begin{remark}[Regularity verification]\label{rem:Xsamp-regularity}
$X^A$ satisfies the regular cell complex conditions of
\cite[\S 2, Def.~1]{HansenGhrist2019}, requiring for each cell of
dimension $n_\alpha$ a homeomorphism of a closed ball in
$\mathbb{R}^{n_\alpha}$ onto the closed cell, mapping the open ball
homeomorphically onto the cell. Each $0$-cell is a single point,
which is the closed unit ball in $\mathbb{R}^0$ and satisfies the
condition trivially. For each $1$-cell $\tau_i$, the closed unit
ball in $\mathbb{R}^1$ is the interval $[0,1]$, whose boundary
$S^0 = \{0, 1\}$ must map to two \emph{distinct} $0$-cells of $X^A$;
the assignment $0 \mapsto \sigma_B$, $1 \mapsto \sigma_\infty$
realizes the required homeomorphism. The double $0$-cell construction
is what verifies \emph{Cond.~4} of \cite[\S 2, Def.~1]{HansenGhrist2019}
(boundaries of $1$-cells must consist of two distinct $0$-cells);
Cond.~(1)--(3) hold automatically by the cell-complex setup of
\cref{def:Xsamp}. Functoriality is vacuous on a $1$-dimensional cell
complex (there are no chains of three distinct cells with strict
inclusions). Hence $X^A$ is a regular cell complex in the sense of
\cite[\S 2, Def.~1, Cond.~4]{HansenGhrist2019}.
\end{remark}

\begin{definition}[Sampling sheaf]\label{def:sampling-sheaf}
The \emph{sampling sheaf} $\samps$ on $X^A$ is the cellular sheaf
defined by
\begin{itemize}[leftmargin=2em]
  \item $\samps(\sigma_B) := V_B$ (with the inner product inherited
        from $\C^N$);
  \item $\samps(\sigma_\infty) := 0$ (the zero Hilbert space);
  \item $\samps(\tau_i) := \C$ for each $i \in V_s$;
  \item the restriction maps are
        $\samps(\sigma_B \le \tau_i) := \eval_i$, where
        $\eval_i : V_B \to \C$, $c \mapsto c_i$, is evaluation at the
        $i$-th index, and
        $\samps(\sigma_\infty \le \tau_i) := 0$, the unique linear
        map from the zero space to $\C$.
\end{itemize}
\end{definition}

\begin{remark}[Cellular-sheaf framework]\label{rem:samps-cellular}
$\samps$ satisfies the cellular sheaf axioms of
\cite[\S 2, Def.~4]{HansenGhrist2019}: each stalk is a
finite-dimensional Hilbert space (the zero space being the zero
object of $\Hilb^{\fin}$, the category of finite-dimensional
Hilbert spaces); each restriction map is a bounded linear operator
between finite-dimensional Hilbert spaces; functoriality is vacuous
since $X^A$ has no chains of three distinct cells with strict
inclusions.
\end{remark}

\begin{lemma}[Canonical isometry $V_B \oplus 0 \cong V_B$]\label{lem:canonical-isometry}
The Hilbert direct sum $V_B \oplus \{0\} = \{(c, 0) : c \in V_B\}$
is canonically isometric to $V_B$ via the projection
$\pi_B(c, 0) := c$, with inverse $\iota_B(c) := (c, 0)$. Both maps
are isometries:
$\lVert (c, 0) \rVert^2_{V_B \oplus 0} = \lVert c \rVert^2_{V_B}$
\cite[\S II.1]{ReedSimonI}.
\end{lemma}

\begin{lemma}[Cochain spaces of $\samps$]\label{lem:samps-cochain-spaces}
For $\samps$ as in \cref{def:sampling-sheaf},
\[
  C^0(\samps) \;=\; V_B \oplus \{0\} \;\cong\; V_B
  \qquad \text{and} \qquad
  C^1(\samps) \;=\; \bigoplus_{i \in V_s} \C \;\cong\; \C^M,
\]
the first isomorphism being the canonical isometry of
\cref{lem:canonical-isometry}.
\end{lemma}

\begin{proof}
$C^0(\samps)$ is the direct sum of the stalks on $0$-cells:
$\samps(\sigma_B) \oplus \samps(\sigma_\infty) = V_B \oplus 0$, and
\cref{lem:canonical-isometry} provides the canonical isometry to
$V_B$. $C^1(\samps)$ is the direct sum of the stalks on $1$-cells,
each a copy of $\C$, indexed by $V_s$.
\end{proof}

\begin{proposition}[Coboundary equals sampling operator]\label{prop:delta-equals-A}\footnote{Formalized in Lean 4 by two statements in \texttt{lean/MRI/Cellular/Sampling.lean}: \texttt{coboundary\_eq\_sampling\_op} proves the operator-layer equality after the canonical map $E \to V_B \oplus 0$, and \texttt{coboundaryHilbert\_eq\_samplingOpEuclidean} proves the Hilbert-layer/PiLp version used for singular values. Full correspondence: \texttt{lean/CROSSREF.md}.}
Under the canonical isometry of
\cref{lem:canonical-isometry,lem:samps-cochain-spaces}, the
coboundary $\delta^0 : C^0(\samps) \to C^1(\samps)$ of $\samps$
agrees with the sampling operator $A : V_B \to \C^M$ of
\cref{eq:A-def}:
\begin{equation}\label{eq:delta0-equals-A}
  \delta^0 \;=\; A \qquad
  \text{as linear maps, after the canonical identification of
        \cref{lem:canonical-isometry,lem:samps-cochain-spaces}.}
\end{equation}
The identification is canonical (uniquely determined by the
sheaf data), so no choice is involved in the comparison.
\end{proposition}

\begin{proof}
By \cref{def:coboundary,def:sampling-sheaf}, for $(c, 0) \in V_B \oplus 0$,
\begin{align*}
  \bigl(\delta^0 (c, 0)\bigr)(\tau_i)
  &= [\sigma_B : \tau_i] \cdot \samps(\sigma_B \le \tau_i)(c)
   + [\sigma_\infty : \tau_i] \cdot \samps(\sigma_\infty \le \tau_i)(0) \\
  &= (+1) \cdot \eval_i(c) + (-1) \cdot 0 \\
  &= c_i
\end{align*}
for each $i \in V_s$. Stacking over $i \in V_s$,
$\delta^0 (c, 0) = (c_{i_1}, \ldots, c_{i_M}) = A c$. Composing with
the canonical isometry of \cref{lem:canonical-isometry} (which is
itself an isometry, hence does not change operator norms or
spectra), $\delta^0 = A$ holds as linear maps after the canonical
identification.
\end{proof}

% =============================================================
\subsection{Strict equality on the sampling sheaf}\label{subsec:sheafbw-strict-equality}
% =============================================================

The reduction of \cref{def:sheafBW} on the sampling sheaf is now a
direct consequence of \cref{prop:delta-equals-A}.

\begin{theorem}[Strict equality]\label{thm:strict-equality}\footnote{Formalized in Lean 4: the paper's full-column-rank form corresponds to \texttt{strict\_equality\_full\_rank} in \texttt{lean/MRI/Cellular/SamplingSpectral.lean} (the equality $\sheafBW(\samps) = \sigma_{\min}(A)/\sigma_{\max}(A)$ with domain-last $\sigma_{\min}$ and the interval $(0,1]$); the upper-boundary iff corresponds to \texttt{sheafBW\_Samps\_eq\_one\_iff\_singularValues\_domain\_eq}. The helper \texttt{strict\_equality} proves the more general positive-spectrum version indexed by \texttt{finrank (range A) - 1}. The further geometric interpretation ``iff $A$ is a positive scalar multiple of an isometry on $V_B$'' is a corollary of $\sigma_{\min}(A) = \sigma_{\max}(A)$ derived in the proof body below and is not separately formalized as a Lean theorem. Full correspondence: \texttt{lean/CROSSREF.md}.}
Let $V$, $V_s$, $B$ be as in \cref{subsec:tri-setup} (and hence
$V_B = V_B(V) \subset \C^N$ and $A : V_B \to \C^M$ as in
\cref{eq:A-def}). Assume $A$ has full column rank, equivalently
$\sigma_{\min}(A) > 0$, equivalently $H^0(\samps) = 0$ (the three
formulations are mutually equivalent by \cref{cor:H0-vanishing}).
For the sampling sheaf $\samps$ on the abstract sampling cell
complex $X^A$ of \cref{def:Xsamp,def:sampling-sheaf},
\begin{equation}\label{eq:strict-equality}
  \sheafBW(\samps) \;=\; \frac{\sigma_{\min}(A)}{\sigma_{\max}(A)}
  \;\in\; (0, 1].
\end{equation}
The upper boundary $\sheafBW(\samps) = 1$ is attained iff
$\sigma_{\min}(A) = \sigma_{\max}(A)$, i.e., iff $A$ is a positive
scalar multiple of an isometry on $V_B$ (cf.\ design clause~(c) of
\cref{def:sheafBW}).
\end{theorem}

\begin{proof}
By \cref{prop:delta-equals-A},
$\delta^0 = A$ as an operator $V_B \to \C^M$ (under the canonical
isometry $C^0(\samps) \cong V_B$). The spectrum of
$(\delta^0)^* \delta^0$ therefore equals the spectrum of $A^* A$,
which is the set of squared singular values of $A$:
$\spec(A^* A) = \{\sigma_l(A)^2\}_{l=0}^{K-1}$. Hence
\begin{align*}
  \gamma_+(\delta^0)
  &= \inf\{\,\lambda \in \spec(A^* A) : \lambda > 0\,\}
   = \sigma_{\min}(A)^2, \\
  \Gamma_+(\delta^0)
  &= \sup \spec(A^* A)
   = \sigma_{\max}(A)^2.
\end{align*}
Substituting into \cref{eq:sheafBW-def},
\[
  \sheafBW(\samps)
  \;=\; \sqrt{\frac{\gamma_+(\delta^0)}{\Gamma_+(\delta^0)}}
  \;=\; \sqrt{\frac{\sigma_{\min}(A)^2}{\sigma_{\max}(A)^2}}
  \;=\; \frac{\sigma_{\min}(A)}{\sigma_{\max}(A)},
\]
each equality being a strict equality of real numbers (no asymptotic
or almost-equal qualification); the operator-level identification
$\delta^0 = A$ used at the start of the proof is canonical
(\cref{prop:delta-equals-A}).
\end{proof}

\begin{corollary}[$H^0$ vanishing on the sampling sheaf]\label{cor:H0-vanishing}\footnote{Formalized in Lean 4: the two statements of this corollary correspond to (i) \texttt{globalSections\_comap\_eq\_ker\_samplingOp} in \texttt{lean/MRI/Cellular/TriangleEquivalence.lean} (the equality $H^0(\samps) = \ker A$ expressed as a \texttt{Submodule.comap} via the canonical isometry); and (ii) \texttt{globalSections\_eq\_bot\_iff\_singularValues\_last\_pos} in the same file (the iff $H^0(\samps) = 0 \iff \sigma_{\min}(A) > 0$). The intermediate \texttt{H0\_vanishing\_iff} in \texttt{lean/MRI/Cellular/SamplingSpectral.lean} bridges to the \texttt{Injective} form. Full correspondence: \texttt{lean/CROSSREF.md}.}
$H^0(\samps) = \ker A$. In particular, $H^0(\samps) = 0$ if and
only if $\sigma_{\min}(A) > 0$.
\end{corollary}

\begin{proof}
$H^0(\samps) = \ker \delta^0$ by the standard cellular-sheaf
identification of $0$-cohomology with global sections; by
\cref{prop:delta-equals-A}, $\ker \delta^0 = \ker A$, and
$\ker A = 0$ iff $\sigma_{\min}(A) > 0$.
\end{proof}

\begin{remark}[Rank-deficient locus]\label{rem:rank-deficient-locus}
When $\sigma_{\min}(A) = 0$ (equivalently $H^0(\samps) \ne 0$), the
hypothesis of \cref{thm:strict-equality} fails and the two sides of
\cref{eq:strict-equality} disagree: the candidate
$\sheafBW^{\mathrm{cand}}(\samps)$ of
\cref{def:sheafbw-candidate} is set to $0$ by convention, while the
intrinsic invariant $\sheafBW(\samps) = \sqrt{\gamma_+(\delta^0) /
\Gamma_+(\delta^0)}$ measures the spectral-gap ratio of $\delta^0$
restricted to $(\ker \delta^0)^\perp$ and remains strictly positive
whenever $\delta^0 \ne 0$ on $(\ker \delta^0)^\perp$. The two
quantities encode different information on the rank-deficient locus:
the candidate records by convention the failure of full sampling,
the intrinsic invariant records the conditioning of $\delta^0$ on
its non-trivial range. \Cref{thm:strict-equality} is the statement
that on the well-conditioned locus $\sigma_{\min}(A) > 0$ the two
quantities coincide strictly.
\end{remark}

\Cref{thm:strict-equality} answers clause~\textup{(R2)} of
\cref{subsec:tri-narrowing} in the well-conditioned regime: the
categorically intrinsic sheaf bandwidth of \cref{def:sheafBW}
reduces to the conditioning ratio
$\sigma_{\min}(A)/\sigma_{\max}(A)$ on the sampling sheaf with
strict equality. \Cref{rem:rank-deficient-locus} delineates the
rank-deficient boundary of the strict equality.

% =============================================================
\subsection{Unitary subcategory and dagger structure}\label{subsec:sheafbw-unitary-subcat}
% =============================================================

\Cref{prop:sheafBW-unitary} required the existence of a unitary
sheaf morphism between $\mathscr{A}$ and $\mathscr{A}'$. We make
the categorical setting explicit and identify its dagger structure.

\begin{definition}[Unitary subcategory]\label{def:unitary-subcat}\footnote{Not formalized in Lean 4: the categorical packaging of the unitary subcategory awaits a fuller cellular-sheaf category in mathlib. The core stalk-wise unitary structure used in the invariance proposition above is captured by \texttt{UnitarySheafMorphism} in \texttt{lean/MRI/Cellular/Unitary.lean}; see \texttt{lean/CROSSREF.md}.}
Let $X$ be a finite regular cell complex. The category
$\CellSh^{\fin}_{\Hilb}(X)$ has finite-dimensional cellular sheaves
of Hilbert spaces on $X$ as objects, and stalk-wise families of
linear operators $\Phi = \{\Phi_\sigma\}_{\sigma \in P_X}$
satisfying the cellular-sheaf compatibility as morphisms. Its
\emph{unitary subcategory}, $\CellSh^{\fin}_{\Hilb,\mathrm{u}}(X)$,
is the subcategory whose morphisms are those $\Phi$ with each
$\Phi_\sigma$ a unitary operator.
\end{definition}

\begin{proposition}[Unitary subcategory is a groupoid]\label{prop:unitary-groupoid}\footnote{Not formalized in Lean 4; the groupoid/dagger structure on the unitary subcategory is left informal pending a fuller categorical formalization in mathlib. See \texttt{lean/CROSSREF.md}.}
The category $\CellSh^{\fin}_{\Hilb,\mathrm{u}}(X)$ is a groupoid.
On this subcategory the dagger functor
$\Phi^\dagger := \{(\Phi_\sigma)^*\}_{\sigma \in P_X}$ coincides
with the inverse functor: $\Phi^\dagger = \Phi^{-1}$.
\end{proposition}

\begin{proof}
Let $\Phi : \mathscr{A} \to \mathscr{A}'$ be unitary, so each
$\Phi_\sigma : \mathscr{A}(\sigma) \to \mathscr{A}'(\sigma)$ is a
unitary operator on a finite-dimensional Hilbert space; in particular
$\Phi_\sigma^{-1} = \Phi_\sigma^*$ \cite[\S VI.1]{ReedSimonI}. To
show $\Phi$ is invertible in $\CellSh^{\fin}_{\Hilb,\mathrm{u}}(X)$,
we verify that the family
$\{\Phi_\sigma^{-1}\}_{\sigma \in P_X}$ defines a sheaf morphism.
The sheaf morphism axiom for $\Phi$ reads
$\Phi_\tau \circ \mathscr{A}(\sigma \le \tau) =
\mathscr{A}'(\sigma \le \tau) \circ \Phi_\sigma$ for every
$\sigma \le \tau$. Applying $\Phi_\tau^{-1}$ on the left and
$\Phi_\sigma^{-1}$ on the right yields
$\mathscr{A}(\sigma \le \tau) \circ \Phi_\sigma^{-1} =
\Phi_\tau^{-1} \circ \mathscr{A}'(\sigma \le \tau)$, the sheaf
morphism axiom for $\{\Phi_\sigma^{-1}\}$ in the opposite direction.
Each $\Phi_\sigma^{-1}$ is unitary (inverse of a unitary is
unitary), so $\Phi^{-1} \in \CellSh^{\fin}_{\Hilb,\mathrm{u}}(X)$.
The identity $\Phi^\dagger_\sigma = \Phi_\sigma^* = \Phi_\sigma^{-1}$
extends stalk-wise to $\Phi^\dagger = \Phi^{-1}$. The category is
therefore a groupoid \cite[\S I.5]{MacLane1971}.
\end{proof}

\begin{remark}[Dagger structure on the unitary subcategory is
trivial]\label{rem:trivial-dagger}
On a groupoid the dagger functor of \cref{prop:unitary-groupoid}
\emph{coincides} with the inverse functor; the dagger axioms of
\cite[\S 2.3]{HeunenVicary2019} are then satisfied by the
groupoid structure alone. This corresponds to the
\emph{dagger-core} of a $\dagger$-category in the standard
nomenclature: it is the maximal sub-groupoid on which dagger and
inverse agree. The substantive content of dagger category
theory---the distinction of dagger from inverse on non-isomorphism
morphisms (isometries, partial isometries, dagger-monics,
dagger-kernels)---is absent on the unitary subcategory by
construction.
\end{remark}

% =============================================================
\subsection{Open Problem: non-trivial dagger on the full category}\label{subsec:sheafbw-op}
% =============================================================

The trivial dagger structure on the unitary subcategory of
\cref{rem:trivial-dagger} leaves open the question of whether a
non-trivial dagger structure exists on the full category
$\CellSh^{\fin}_{\Hilb}(X)$.

\begin{openproblem}[Non-trivial dagger on $\CellSh^{\fin}_{\Hilb}(X)$]
\label{op:non-trivial-dagger}
Does there exist a contravariant involutive functor
\[
  \dagger \;:\; \CellSh^{\fin}_{\Hilb}(X)
  \longrightarrow \CellSh^{\fin}_{\Hilb}(X)^{\mathrm{op}}
\]
satisfying the standard dagger-functor axioms of
\cite[\S 2.3, 2.18--2.19]{HeunenVicary2019}, namely
\begin{enumerate}[label=(D\arabic*)]
  \item $\mathscr{A}^\dagger = \mathscr{A}$ for every object;
  \item $(\mathrm{id}_{\mathscr{A}})^\dagger = \mathrm{id}_{\mathscr{A}}$;
  \item $(\Phi \circ \Psi)^\dagger = \Psi^\dagger \circ \Phi^\dagger$;
  \item $(\Phi^\dagger)^\dagger = \Phi$;
\end{enumerate}
that restricts on the unitary subcategory
$\CellSh^{\fin}_{\Hilb,\mathrm{u}}(X)$ to the stalk-wise Hilbert
adjoint of \cref{prop:unitary-groupoid}, and is moreover
\emph{non-trivial} on non-unitary morphisms in the sense of
distinguishing isometries from dagger-monics from arbitrary morphisms?
\end{openproblem}

The obstruction to a stalk-wise definition is recorded for context.
On a non-unitary morphism $\Phi$, the stalk-wise Hilbert adjoints
$\{(\Phi_\sigma)^*\}_{\sigma \in P_X}$ in general fail the cellular
sheaf morphism axiom: starting from
$\Phi_\tau \circ \mathscr{A}(\sigma \le \tau) =
\mathscr{A}'(\sigma \le \tau) \circ \Phi_\sigma$ and taking adjoints
yields
$\mathscr{A}(\sigma \le \tau)^* \circ \Phi_\tau^* =
\Phi_\sigma^* \circ \mathscr{A}'(\sigma \le \tau)^*$, which is the
sheaf morphism axiom for the \emph{adjoint restriction maps}
$\mathscr{A}(\sigma \le \tau)^*, \mathscr{A}'(\sigma \le \tau)^*$
running in the opposite direction in the face poset. Without
unitarity of restriction maps---in general absent in cellular
sheaves of Hilbert spaces---the stalk-wise dagger fails to be a
sheaf morphism. A non-trivial dagger functor on the full category
must therefore use a non-stalk-wise extension, with no precedent in
the cellular-sheaf literature: \cite{LahtinenStenvall2020} treats
dagger structure on the category of finite element discretizations
(distinct from cellular sheaves), and \cite{HansenGhrist2019}
introduces the unitary subcategory but does not extend a dagger
beyond it.

\Cref{op:non-trivial-dagger} is independent of the main results of
this paper. \Cref{prop:sheafBW-unitary,thm:strict-equality} hold
under unitary equivalence in
$\CellSh^{\fin}_{\Hilb,\mathrm{u}}(X)$, which is sufficient for the
intrinsicness clause~\textup{(R1)} of \cref{subsec:tri-narrowing}
in the form proven here. A non-trivial dagger on the full category
would strengthen \cref{prop:sheafBW-unitary} to dagger invariance,
but is not required for either the strict equality of
\cref{thm:strict-equality} or for the three-instance instantiation
of \cref{sec:three_instances}.

% --- Closing hook (to Section 6) ---
\Cref{def:sheafBW,prop:sheafBW-unitary,thm:strict-equality} answer
clauses~\textup{(R1)} and~\textup{(R2)} of
\cref{subsec:tri-narrowing}: the sheaf bandwidth is defined
intrinsically from the cellular-sheaf data, is invariant under
unitary equivalence, and reduces to $\sigma_{\min}(A)/\sigma_{\max}(A)$
on the sampling sheaf $\samps$ with strict equality. To exhibit
the framework's reach beyond the canonical sampling instance, we now
apply it to three previously disjoint domains---signal processing on
simplicial complexes, finite element electromagnetics, and Brillouin
zone tetrahedron quadrature---unified by a common categorical template.


% =============================================================
% Section 6. Three Instances
% =============================================================
\section{Three Instances: TopSP, Finite Element EM, Brillouin Zone}\label{sec:three_instances}
% =============================================================
% Section 6. Three Instances: TopSP, Finite Element EM, Brillouin Zone
%
% Source: A-route-ch3-three-instances-v0.2.md
% Template: paper/v0.1/notes/narrative.md §4 §6
% Cross-refs to §5: \cref{def:sheafBW}, \cref{prop:sheafBW-unitary},
%                    \cref{thm:strict-equality}, \cref{def:Xsamp},
%                    \cref{def:sampling-sheaf}, \cref{op:non-trivial-dagger}
% =============================================================

\noindent
We instantiate the framework of \cref{sec:sheafbw} in three previously
disjoint domains. Each instance specializes the abstract domain $V_?$
to a domain-specific Hilbert subspace and the abstract operator $A$ to
a domain-specific evaluation or quadrature operator, while the
categorical structure---the abstract sampling cell complex $X^A$
(\cref{def:Xsamp}), the sampling sheaf $\samps$
(\cref{def:sampling-sheaf}), the strict equality
\cref{thm:strict-equality}, and the unitary invariance
\cref{prop:sheafBW-unitary}---is shared verbatim across the three
instances. The categorical structure is unified verbatim, but the
\emph{physical interpretations} differ instance by instance: in
particular, the sampling map $S$ and integration map $I$ introduced in
\cref{subsec:6-maps} are formal scalar maps on the geometry data $V$,
and the role they play (evaluation versus quadrature, real versus
synthetic data, independence versus coupling of $S$ and $I$) is
instance-specific. The \emph{verbatim} clause refers to the categorical
structure only, not to the physical correspondence; the explicit
divergences are recorded in \cref{tab:6-instances} and in the
per-instance discussion of \cref{subsec:6-bz}.

% -------------------------------------------------------------
\subsection{Common setup}\label{subsec:6-common}

Throughout this section we fix a finite-dimensional Hilbert subspace
$V_? \subset \C^N$ of dimension $K$, a sampling location set
$V_s \subset \{1, \dots, N\}$ of size $M$, and the corresponding
sampling operator
\[
  A : V_? \to \C^M,
  \qquad
  c \mapsto \bigl(c_{i_n}\bigr)_{n=1}^{M},
\]
where $\{i_n\}_{n=1}^{M}$ enumerates $V_s$. The abstract sampling cell
complex $X^A$ is defined by \cref{def:Xsamp} (one principal $0$-cell
$\sigma_B$, an auxiliary $0$-cell $\sigma_\infty$, and $M$ $1$-cells
$\{\tau_i\}_{i \in V_s}$). The sampling sheaf $\samps$ is defined by
\cref{def:sampling-sheaf} with stalks $\samps(\sigma_B) = V_?$,
$\samps(\sigma_\infty) = 0$, $\samps(\tau_i) = \C$, and restriction
maps $\eval_i$ at $\sigma_B$ and zero at $\sigma_\infty$.

By \cref{thm:strict-equality}, under the hypothesis
$\sigma_{\min}(A) > 0$ (equivalently $H^0(\samps_{V_s, V_?}) = 0$),
\begin{equation}\label{eq:6-strict}
  \sheafBW(\samps_{V_s, V_?})
  \;=\;
  \frac{\sigma_{\min}(A)}{\sigma_{\max}(A)},
\end{equation}
and by \cref{prop:sheafBW-unitary} this quantity is invariant under
unitary equivalence of cellular sheaves over $X^A$. Each of the three
instances below verifies the full-column-rank hypothesis under its
standard configuration: Instance~1 by direct numerical check on
$32$ MRI configurations (\cref{subsec:6-topsp}), Instance~2 by the
$H^1$-Rayleigh basis of $V_S^{\FEM}$ being linearly independent at
the chosen vertex sample (\cref{subsec:6-femem}), and Instance~3 by
non-degeneracy of the linear tetrahedron weight matrix
$w^{\mathrm{lin}}$ on the occupied subspace
(\cref{subsec:6-bz}).

The three instances differ only in the choice of $(V_?, A)$. We display
them in a uniform format and summarize the comparison in
\cref{tab:6-instances}.

% -------------------------------------------------------------
\subsection{Instance 1: Topological signal processing}\label{subsec:6-topsp}

\paragraph{Domain $V_B^{\TopSP}$.}
Let $X(V)$ be a cellular complex with vertex set
$V = \{k_1, \dots, k_N\} \subset \Omega_k$, and let $L_0 = D - W$ be its
graph Laplacian with eigendecomposition
$(\mu_l, v_l)_{l=0}^{N-1}$, $\mu_0 \le \dots \le \mu_{N-1}$. Following
\cite[\S 4.4]{BarbarossaSardellitti2020}, the Hodge bandlimited
subspace at threshold $B$ is
\[
  V_B^{\TopSP}
  \;=\;
  \operatorname{span}\{\, v_l : \mu_l \le B \,\}
  \;\subset\; \C^N,
  \qquad
  K \;=\; \dim V_B^{\TopSP}.
\]

\paragraph{Operator $A^{\TopSP}$.}
With sampling location set $V_s$ of size $M$, the sampling operator is
\[
  A^{\TopSP}
  \;:\;
  V_B^{\TopSP} \to \C^M,
  \qquad
  A^{\TopSP}_{n,l}
  \;=\;
  v_l(i_n),
  \qquad n = 1, \dots, M,\;\; l = 0, \dots, K-1.
\]

\paragraph{Sheaf bandwidth.}
By \cref{eq:6-strict},
\[
  \sheafBW(\samps_{V_s, V_B^{\TopSP}})
  \;=\;
  \frac{\sigma_{\min}(A^{\TopSP})}{\sigma_{\max}(A^{\TopSP})}.
\]

\paragraph{Empirical evidence on simulated MRI sampling geometries.}
A simulation suite developed in the present project applied
$\sigma_{\min}(A^{\TopSP})/\sigma_{\max}(A^{\TopSP})$ to $32$
configurations covering five trajectory families (Cartesian, random,
sparkling, spiral, radial) over up to seven subsampling levels per
family ($M$ ranging from $52$ to $305$, at base scale $N$ between
$288$ and $305$; the central-symmetric inner-half level exists only
for the spiral and radial families, giving $32$ rather than the full
$5 \times 7$ grid). All $32$ configurations attained full column rank,
with Cartesian ordering above radial at every matched subsampling
level; at the most rank-stressed level the separation reaches more
than two orders of magnitude (Cartesian $\sheafBW \approx 0.0091$
versus radial $\sheafBW \approx 2.07 \times 10^{-5}$). This ordering is
consistent with the reconstruction quality reported in MRI sampling
practice \cite{LustigDonohoPauly2007,LazarusEtAl2019Sparkling}. The
numerical study is reported in \cref{sec:numerical}.

% -------------------------------------------------------------
\subsection{Instance 2: Finite element electromagnetics}\label{subsec:6-femem}

\paragraph{Setup.}
Let $\Omega \subset \R^3$ be a bounded Lipschitz domain with a
simplicial mesh $K(\tau_h)$ (each top-dimensional simplex a
tetrahedron). Following \cite{Whitney1957,Bossavit1988}, the Whitney
$0$-form space is
\[
  W^0(K)
  \;=\;
  \operatorname{span}\{\, \lambda_i : i = 1, \dots, N \,\}
  \;\cong\; \C^N
  \qquad
  (\lambda_i \text{ the hat function at vertex } x_i).
\]

\paragraph{Domain $V_S^{\FEM}$.}
Let $\{w_l\}_{l=0}^{K-1}$ be the lowest $K$ eigenmodes of the
$H^1$-Rayleigh quotient on $W^0(K)$, normalized so that
$\langle w_l, w_m \rangle_{L^2} = \delta_{lm}$:
\[
  \langle \nabla w_l, \nabla \varphi \rangle_{L^2}
  \;=\;
  \lambda_l \, \langle w_l, \varphi \rangle_{L^2}
  \quad \forall \varphi \in W^0(K),
  \qquad
  \lambda_0 \le \dots \le \lambda_{K-1}.
\]
Define
\[
  V_S^{\FEM}
  \;:=\;
  \operatorname{span}\{\, w_l : l = 0, \dots, K-1 \,\}
  \;\subset\; \C^N.
\]
Then $V_S^{\FEM}$ is a $K$-dimensional Hilbert subspace of $\C^N$.

\paragraph{Operator $A^{\FEM}$.}
Choose $M$ vertices $\{x_{i_n}\}_{n=1}^{M}$ as the sampling locations.
The sampling operator is
\[
  A^{\FEM} : V_S^{\FEM} \to \C^M,
  \qquad
  (A^{\FEM} c)_n
  \;=\;
  \sum_{l=0}^{K-1} c_l \, w_l(x_{i_n}),
\]
i.e.\ $A^{\FEM}_{n,l} = w_l(i_n)$, the same matrix form as
$A^{\TopSP}$ with the basis $\{v_l\}$ replaced by the
$H^1$-Rayleigh modes $\{w_l\}$. The pair $(V_S^{\FEM}, A^{\FEM})$
satisfies the input requirements of \cref{thm:strict-equality}: (i)
$V_S^{\FEM}$ is a finite-dimensional Hilbert subspace by
construction; (ii) $A^{\FEM}$ is bounded and linear, automatic in
finite dimensions; (iii) $\sigma_{\min}(A^{\FEM}) > 0$ holds under
the standard non-degeneracy condition that the $H^1$-Rayleigh modes
$\{w_l\}_{l=0}^{K-1}$ are linearly independent on the chosen vertex
sample $\{x_{i_n}\}_{n=1}^{M}$, a mild assumption on the mesh and
sample configuration.

\paragraph{Project-internal scope: relation to AFW bounded cochain projection.}
\cite{AFW2010} construct on the continuous Hilbert complex
$(W^k(\Omega), d^k)$ a bounded cochain projection $\pi_h$ satisfying
$\pi_h \circ \iota = \mathrm{id}_{V_h}$ and $d \circ \pi_h = \pi_h \circ
d$ (\cite[\S 3.1.1]{AFW2010}). The operator $A^{\FEM}$ above and the
AFW projector $\pi_h$ are distinct objects:

\begin{center}
\begin{tabular}{lll}
\toprule
& $A^{\FEM}$ (this paper) & $\pi_h$ (\cite{AFW2010}) \\
\midrule
Domain        & $V_S^{\FEM} \subset \C^N$ & $H \Lambda^k(\Omega)$ \\
Codomain      & $\C^M$ (vertex evaluations) & finite-dim $V_h$ \\
Role          & sampling operator & projection / interpolation \\
\bottomrule
\end{tabular}
\end{center}

The hypothesis of \cref{thm:strict-equality} requires a
finite-dimensional Hilbert subspace and a bounded linear operator; it
does not require a bounded cochain projection. Instance~2 is
established within the finite-dimensional restriction of
\cref{subsec:sheafbw-hilbert-complex}; the continuous-side extension to
$H^1(\Omega)$ via $\pi_h$ is outside the scope of the present paper.

\paragraph{Empirical evidence.}
Instance~2 is established as a \emph{scope-bounded categorical claim}:
the construction lies within the finite-dimensional restriction
discussed above, and the symbolic identification $A^{\FEM} \cong \pi_h$
holds verbatim under that restriction. No independent empirical channel
(numerical study or established practice) is provided for Instance~2 in
this paper; the evidence channels listed in \cref{tab:6-evidence} record
this explicitly. The numerical verification in \cref{sec:numerical}
covers the algebraic identity \cref{eq:6-strict} on Instance~1 and does
not probe Instance~2 separately. The continuous-side extension to
$H^1(\Omega)$ via the bounded cochain projection $\pi_h$ is outside the
present paper's scope and is recorded as an engineering outlook in
\cref{sec:discussion}.

\paragraph{Sheaf bandwidth.}
By \cref{eq:6-strict},
\[
  \sheafBW(\samps_{V_s, V_S^{\FEM}})
  \;=\;
  \frac{\sigma_{\min}(A^{\FEM})}{\sigma_{\max}(A^{\FEM})}.
\]

% -------------------------------------------------------------
\subsection{Instance 3: Brillouin zone tetrahedron quadrature}\label{subsec:6-bz}

\paragraph{Setup.}
Let $\Omega_{\BZ} \subset \R^3$ be the first Brillouin zone of a
crystalline solid, equipped with a tetrahedron tessellation by
reciprocal-space points $\{\mathbf{k}_i\}_{i=1}^{N}$. For a fixed band
index, let $\epsilon_i$ be the band energy at $\mathbf{k}_i$ and
$\epsilon_F$ the Fermi level.

\paragraph{Domain $V_{\occ}^{\BZ}$.}
In the linear-segment regime of the Bl\"ochl tetrahedron method
\cite[\S II.B]{Blochl1994}, the band energy $\epsilon(\mathbf{k})$ is
linearly interpolated within each tetrahedron and the Heaviside step
$\theta(\epsilon_F - \epsilon)$ is integrated exactly over this linear
segment, yielding the linear tetrahedron weights $w^{\mathrm{lin}}_{i_n, j}$
of \cite{LehmannTaut1972,JepsenAndersen1971,Blochl1994}. We define
\[
  V_{\occ}^{\BZ}
  \;:=\;
  \operatorname{span}\{\, \chi_i : i \in \mathrm{occ}(\epsilon_F) \,\}
  \;\subset\;
  \C^N
\]
as the span of Whitney $0$-form vertex hat functions $\chi_i$ at
$\mathbf{k}_i$ over occupied $k$-points
$\mathrm{occ}(\epsilon_F) = \{i : \epsilon_i \le \epsilon_F\}$;
$K = |\mathrm{occ}(\epsilon_F)|$.

\paragraph{Scope of $V_{\occ}^{\BZ}$.}
The construction of $V_{\occ}^{\BZ}$ as a Whitney-$0$-form span is a
project-internal categorical reformulation of the linear-tetrahedron
framework; it is \emph{not} a standard ``occupied subspace'' of band
theory (the standard object is the span of Bloch eigenfunctions of
occupied bands, residing in a distinct vector space). The
reformulation is motivated by the fact that the linear tetrahedron
weights $w^{\mathrm{lin}}_{i_n, j}$ of
\cite{Blochl1994,LehmannTaut1972,JepsenAndersen1971} are defined on
this hat-function basis, making $V_{\occ}^{\BZ}$ the natural domain
for the operator $A^{\BZ}$ of the next paragraph.

\paragraph{Operator $A^{\BZ}$.}
With $M$ test points (a sub-tetrahedral integration grid) at vertex
locations,
\[
  A^{\BZ} : V_{\occ}^{\BZ} \to \C^M,
  \qquad
  A^{\BZ}_{n, j}
  \;=\;
  w_{i_n j}^{\mathrm{lin}},
\]
where $w_{i_n j}^{\mathrm{lin}}$ is the linear tetrahedron weight of
\cite{LehmannTaut1972,JepsenAndersen1971}.

\paragraph{Physical role of $A^{\BZ}$ versus $A^{\TopSP}$ and $A^{\FEM}$.}
Across the three instances, the abstract operator $A : V_? \to \C^M$
of \cref{subsec:6-common} plays structurally different physical
roles. In \cref{subsec:6-topsp,subsec:6-femem} the matrix $A$ is a
\emph{row-sampling} operator at vertex locations:
$A^{\TopSP}_{n,l} = v_l(i_n)$ and $A^{\FEM}_{n,l} = w_l(i_n)$ are
bare evaluations of basis functions. In Instance~3 the matrix
$A^{\BZ}_{n,j} = w_{i_n j}^{\mathrm{lin}}$ is the
linear-tetrahedron \emph{integration weight} matrix: the $n$-th row
records the weights $w_{i_n j}^{\mathrm{lin}}$ that the
sub-tetrahedral integration grid attaches to the occupied vertex
basis at test point $i_n$. The two roles---row-sampling and
quadrature weighting---are physically distinct.

The categorical framework of \cref{sec:sheafbw}, however, requires
only that $A$ be a bounded linear map $V_? \to \C^M$; it does not
distinguish sampling-style from quadrature-style operators.
Consequently:
\begin{itemize}[leftmargin=*]
\item the strict equality \cref{eq:6-strict} applies to $A^{\BZ}$
  verbatim, with $\sheafBW$ measuring the conditioning of the
  integration weight matrix on $V_{\occ}^{\BZ}$;
\item the sampling map $S$ and the integration map $I$
  defined in \cref{def:sampling-map,def:integration-map}
  are formal scalar maps on the geometry parameter $V$ and do not
  depend on the physical role of $A$ within an instance. In
  Instance~3, $S^{\BZ}$ and $I^{\BZ}$ both involve the integration
  weight matrix $A^{\BZ}$ and are therefore not physically
  independent within the BZ context, although they remain formally
  distinct as scalar maps of $V$. The directional content of
  \cref{empirical:8b} is established across instances on the
  $V$-parameter side, not within Instance~3 in isolation.
\end{itemize}

\paragraph{Project-internal scope: Bl\"ochl correction.}
\cite{Blochl1994} introduces a quadratic correction term that improves
the tetrahedron quadrature beyond the linear segment. This correction
is a Richardson-extrapolation-style refinement that does not lie in
the Whitney $0$-form framework (\cite[\S II.C]{Blochl1994}).
Instance~3 is established within the linear-segment regime; extending
sheaf bandwidth analysis to the Bl\"ochl-corrected quadrature is
outside the scope of the present paper.

\paragraph{Sheaf bandwidth.}
By \cref{eq:6-strict},
\[
  \sheafBW(\samps_{V_s, V_{\occ}^{\BZ}})
  \;=\;
  \frac{\sigma_{\min}(A^{\BZ})}{\sigma_{\max}(A^{\BZ})}.
\]

\paragraph{Empirical evidence: five decades of practice.}
Production density-functional codes have, since 1971, used
non-Cartesian tessellations
(Monkhorst--Pack \cite{MonkhorstPack1976}; tetrahedron methods
\cite{LehmannTaut1972,JepsenAndersen1971,Blochl1994}) for Brillouin
zone integration, while Cartesian quadrature is documented to perform
worst on smooth band-structure integrands. This empirical hierarchy
is consistent with the Pareto trade-off articulated below
(\cref{empirical:8b}).

% -------------------------------------------------------------
\subsection{Three-instance comparison}\label{subsec:6-table}

\begin{table}[h]
\centering
\begin{tabular}{lccc}
\toprule
& \textbf{Instance 1: TopSP} & \textbf{Instance 2: FEM--EM} & \textbf{Instance 3: BZ} \\
\midrule
$V_?$              & $V_B^{\TopSP}$ (Hodge band)     & $V_S^{\FEM}$ ($H^1$-Rayleigh)   & $V_{\occ}^{\BZ}$ (Whitney $0$-form, occupied) \\
$A_{n,l}$          & $v_l(i_n)$                       & $w_l(i_n)$                        & $w_{i_n j}^{\mathrm{lin}}$ \\
Physical role of $A$ & row-sampling                   & row-sampling                      & integration weights \\
$\sheafBW$         & $\sigma_{\min}/\sigma_{\max}$    & $\sigma_{\min}/\sigma_{\max}$     & $\sigma_{\min}/\sigma_{\max}$ \\
Empirical evidence & $32/32$ MRI configs              & finite-dim restriction            & 50-year practice \\
Reference          & \cite{BarbarossaSardellitti2020} & \cite{AFW2010,Whitney1957}        & \cite{Blochl1994,LehmannTaut1972} \\
\bottomrule
\end{tabular}
\caption{Three instances of the framework of \cref{sec:sheafbw}. The
  categorical structure $X^A + \samps + \cref{thm:strict-equality} +
  \cref{prop:sheafBW-unitary}$ is shared verbatim. The instances differ
  only in the domain-specific choice of $(V_?, A)$; the spectral ratio
  $\sigma_{\min}(A)/\sigma_{\max}(A)$ is the unified invariant.}
\label{tab:6-instances}
\end{table}

% -------------------------------------------------------------
\subsection{Sampling and integration maps}\label{subsec:6-maps}

The unified description suggests two scalar maps on the same data.

\begin{definition}[Sampling map]\label{def:sampling-map}\footnote{Not formalized in Lean 4. The sampling map $S$ is a physical correspondence (B-spline FE, quadrature, Bloch--Floquet $\psi_\theta$, etc.) from geometry $V$ to spectral data, not a mathematical theorem statement. See \texttt{lean/CROSSREF.md} \S C.}
For a fixed $V_?$, the \emph{sampling map} $S$ assigns to a sampling
cell complex $X^A_V$ (parameterized by the geometry $V$) the spectral
ratio
\[
  S(X^A_V)
  \;:=\;
  \sheafBW(\samps_V)
  \;=\;
  \frac{\sigma_{\min}(A_V)}{\sigma_{\max}(A_V)}.
\]
By \cref{prop:sheafBW-unitary}, $S$ is invariant under unitary
equivalence of cellular sheaves over $X^A$.
\end{definition}

\begin{definition}[Integration map]\label{def:integration-map}\footnote{Not formalized in Lean 4. The integration map $I$ is a physical correspondence (quadrature error norm) from geometry $V$ to scalar error, not a mathematical theorem statement. See \texttt{lean/CROSSREF.md} \S C.}
For a fixed $V_?$, the \emph{integration map} $I$ assigns to $X^A_V$
the quadrature error norm
\[
  I(X^A_V)
  \;:=\;
  \biggl\|\, \mathcal{I}^* \;-\; \sum_{i \in V_s} w_i \, \eval_i^*\, \biggr\|_{V_? \to \C},
\]
where $\mathcal{I}^* : V_? \to \C$ is an exact-integration dual
functional and $\{w_i\}_{i \in V_s}$ are quadrature weights, both
chosen instance-specifically (e.g., the linear-tetrahedron weight
$w^{\mathrm{lin}}$ of \cref{subsec:6-bz} for Instance~3).
\end{definition}

The choice of $\mathcal{I}^*$ and $\{w_i\}$ is fixed per instance and
not unique across instances; the across-instance directional content
of \cref{empirical:8b} does not depend on a specific choice (cf.\
the disclaim in \cref{subsec:6-bz} on the formal nature of $S, I$).

% -------------------------------------------------------------
\subsection{Empirical Proposition: $S$--$I$ Pareto trade-off}\label{subsec:6-empirical}

\begin{empirical}[Sampling and integration trace a Pareto frontier under $V$-optimization]\label{empirical:8b}\footnote{Not formalized in Lean 4. This is an \emph{empirical} proposition supported by two quantitative evidence channels plus one heuristic skeleton (see body); it is not promoted to theorem status per the double-track verification principle. See \texttt{lean/CROSSREF.md} \S C.}
Fix the type of $V_?$ (one of $V_B^{\TopSP}$, $V_S^{\FEM}$,
$V_{\occ}^{\BZ}$) and the strategy for choosing $V_s$. As the geometry
$V$ varies over admissible cellular complexes, no $V$ simultaneously
maximizes $V \mapsto S(X^A_V)$ and minimizes $V \mapsto I(X^A_V)$;
the pair $(S, I)$ traces a Pareto frontier in $V$.
\end{empirical}

\paragraph{Heuristic skeleton.}
The strict equality~\cref{eq:6-strict} gives the $S$-side directly.
For the $I$-side, the spectral framework of \cite{Pilleboue2015} on
Euclidean point sets suggests, for sufficiently smooth integrands, a
directional decomposition: the integration error scales with the
high-frequency tail of the integrand expanded in the $V_?$ eigenbasis,
weighted by the Fourier spectrum of the sample-set indicator. This is
a heuristic of directional value only, not a quantitative bound on
cellular-sheaf integration. The directional content of
\cref{empirical:8b} is the following.
\begin{itemize}[leftmargin=*]
\item $S$ is large iff the columns of $A_V$ are close to orthonormal
  on $V_s$, which holds when the geometry $V_s$ is locally coherent
  with the $V_?$-modes $\{v_l\}$. Regular structures (Cartesian grids)
  achieve this local coherence.
\item $I$ is small iff the spectrum of $V_s$ flattens at high
  frequencies, which requires the geometry $V_s$ to break regular
  structure. Blue-noise tessellations \cite{HeitzBelcour2019} and
  Monkhorst--Pack grids \cite{MonkhorstPack1976} exemplify this.
\end{itemize}
The two requirements pull $V$ in opposite directions: sampling favors
regular structure; integration favors broken regularity. This is the
heuristic content of \cref{empirical:8b}.

\paragraph{Evidence and heuristic skeleton.}
The directional content of \cref{empirical:8b} is supported by two
empirical channels and a heuristic skeleton, recorded in
\cref{tab:6-evidence}.

\begin{table}[h]
\centering
\begin{tabular}{llll}
\toprule
\textbf{Type} & \textbf{Item} & \textbf{Source} & \textbf{Strength} \\
\midrule
Empirical & b3pre numerical study ($32$ configs, multi-$N$ robustness) & this paper, \cref{sec:numerical}      & strong \\
Empirical & Brillouin zone practice (five decades)                       & \cite{Blochl1994,LehmannTaut1972,JepsenAndersen1971} & strong \\
Heuristic & Euclidean spectral framework (Fourier-spectrum form)       & \cite{Pilleboue2015}                   & directional only \\
\bottomrule
\end{tabular}
\caption{Two empirical channels and one heuristic skeleton
  supporting \cref{empirical:8b}. The b3pre numerical study and the
  multi-$N$ robustness scan are stratifications of a single
  project-internal study, counted as one empirical channel. The
  Brillouin zone practice channel is independent published evidence
  on $V$-geometry choice for quadrature, not on the joint
  sampling-vs-integration trade-off; it supports the directional
  content via the integration side. The Pilleboue spectral framework
  is a heuristic skeleton on Euclidean Monte Carlo variance, not
  direct evidence on cellular-sheaf integration on simplicial or
  tetrahedral complexes.}
\label{tab:6-evidence}
\end{table}

\paragraph{Status.}
\cref{empirical:8b} is a directional statement supported by two
empirical channels and a heuristic skeleton (\cref{tab:6-evidence});
it is not a constructive Pareto-optimality theorem. Constructive proofs typically
require compactness arguments and convex-optimization duality applied
to the joint $(S, I)$ map. The constructive proof is treated as an
independent open problem (\cref{op:8b-rigor} below). Should that open
problem be resolved, \cref{empirical:8b} upgrades to a theorem.

% -------------------------------------------------------------
\subsection{Open Problem: constructive Pareto proof}\label{subsec:6-op}

\begin{openproblem}[Constructive Pareto proof for \cref{empirical:8b}]\label{op:8b-rigor}
Establish, by compactness arguments and convex-optimization duality on
the joint map $V \mapsto (S(X^A_V), I(X^A_V))$, the formal statement
that no admissible geometry $V$ simultaneously achieves
$S(X^A_V) > S^* - \epsilon$ and $I(X^A_V) < I^* + \epsilon$ for
sufficiently small $\epsilon > 0$, where $(S^*, I^*)$ are the Pareto
extrema. The expected mode of attack:
\begin{enumerate}[leftmargin=*]
\item compactness of the $V$-parameter space (under suitable boundary
  conditions);
\item convex relaxation of the joint map, with the dual optimization
  encoding the trade-off;
\item closing the duality gap via spectral arguments paralleling
  \cite{Pilleboue2015}.
\end{enumerate}
\cref{op:8b-rigor} is independent of the main theorems of this paper:
\cref{empirical:8b} is not invoked in the proof of
\cref{thm:strict-equality}, and the categorical bridge stands without
it.
\end{openproblem}

\noindent
The unified framework reduces all three instances to a spectral ratio
on $A$. We verify this strict equality numerically across the three
instances in \cref{sec:numerical}.


% =============================================================
% Section 7. Numerical Verification
% =============================================================
\section{Numerical Verification}\label{sec:numerical}
% =============================================================
% Section 7. Numerical Verification
%
% Source: A-route-ch3-three-instances-v0.2.md §5 + code/A_route_ch3_numerical.py
% Template: paper/v0.1/notes/narrative.md §4 §7
% =============================================================

\noindent
\cref{thm:strict-equality} reduces $\sheafBW(\samps)$ to a spectral
ratio on the sampling operator $A$. We verify this strict equality
numerically across the three instances of \cref{sec:three_instances}.
The verification confirms the linear-algebraic identity and the code
implementation; it does not constitute independent physical evidence
for each instance, as we make explicit in
\cref{subsec:7-honest-statement}.

% -------------------------------------------------------------
\subsection{Experimental setup}\label{subsec:7-setup}

A self-contained Python script (\texttt{A\_route\_ch3\_numerical.py},
$147$~lines, depending only on
\texttt{numpy} \cite{HarrisEtAl2020NumPy}) constructs a toy operator
$A \in \C^{M \times K}$ for each of the three instances and computes
two quantities for comparison:
\begin{enumerate}[leftmargin=*,label=(\roman*)]
\item the spectral-gap form
  $\sqrt{\gamma_+(\delta^0)/\Gamma_+(\delta^0)}$, where
  $(\delta^0)^* \delta^0 = A^* A$ on the relevant $V_?$ subspace;
\item the singular-value form
  $\sigma_{\min}(A)/\sigma_{\max}(A)$.
\end{enumerate}
For each instance, $A$ is generated by drawing a random orthonormal
basis $U \in \C^{N \times K}$ via QR factorization of a complex
Gaussian matrix and then sub-sampling $M$ rows of $U$ to obtain
$A$, mirroring the row-sampled basis representation
$A_{n,l} = v_l(i_n)$ used uniformly across the three instances of
\cref{sec:three_instances}. The dimensions are
\begin{center}
\begin{tabular}{lcccl}
\toprule
Instance & $N$ & $K$ & $M$ & Subspace label \\
\midrule
1: TopSP & $100$ & $10$ & $30$ & $V_B^{\TopSP}$ \\
2: FEM--EM & $50$ & $6$ & $8$ & $V_S^{\FEM}$ \\
3: BZ & $30$ & $4$ & $12$ & $V_{\occ}^{\BZ}$ \\
\bottomrule
\end{tabular}
\end{center}
The random seed is fixed at $\mathtt{seed} = 20260507$ for
reproducibility. The relative deviation between the two forms is
defined as
\[
  \mathrm{reldiff}
  \;:=\;
  \frac{\bigl|\sqrt{\gamma_+/\Gamma_+} - \sigma_{\min}/\sigma_{\max}\bigr|}{\sigma_{\min}/\sigma_{\max}}.
\]

% -------------------------------------------------------------
\subsection{Results}\label{subsec:7-results}

\begin{table}[h]
\centering
\begin{tabular}{lccc}
\toprule
\textbf{Instance} & $\sqrt{\gamma_+/\Gamma_+}$ & $\sigma_{\min}/\sigma_{\max}$ & $\mathrm{reldiff}$ \\
\midrule
1: TopSP   & $0.367756144311$ & $0.367756144311$ & $7.55 \times 10^{-16}$ \\
2: FEM--EM & $0.053301378959$ & $0.053301378959$ & $1.64 \times 10^{-14}$ \\
3: BZ      & $0.519194438983$ & $0.519194438983$ & $2.14 \times 10^{-16}$ \\
\bottomrule
\end{tabular}
\caption{Numerical verification of \cref{thm:strict-equality} on the
  three instances at $\mathtt{seed} = 20260507$. The two forms
  (spectral-gap and singular-value) of \cref{def:sheafBW} agree to at
  least twelve decimal digits in every instance. Instances~1 and~3
  attain $\mathrm{reldiff}$ within $4 \times \epsilon_{\mathrm{mach}}$
  (IEEE~754 unit roundoff $\epsilon_{\mathrm{mach}} \approx
  2.22 \times 10^{-16}$); Instance~2 deviates by $1.64 \times 10^{-14}$,
  reflecting the expected roundoff amplification of comparing
  $\sqrt{\lambda_{\min}(A^* A)/\lambda_{\max}(A^* A)}$ with
  $\sigma_{\min}(A)/\sigma_{\max}(A)$ when $A$ is poorly conditioned,
  since $\kappa(A^* A) = \kappa(A)^2$. All three deviations are at
  least an order of magnitude below the floating-point error budget for
  this comparison on matrices of the given shapes.}
\label{tab:7-results}
\end{table}

The maximum relative deviation across the three instances is
$1.64 \times 10^{-14}$ (Instance~2), driven by the squaring of the
condition number under the $A^* A$ form rather than by any numerical
failure of \cref{thm:strict-equality}. The spectral-gap form and the
singular-value form agree to within the floating-point error budget
in every instance.

% -------------------------------------------------------------
\subsection{Honest statement on the nature of this verification}\label{subsec:7-honest-statement}

The script of \cref{subsec:7-setup} verifies the spectral-gap form
$\sqrt{\gamma_+/\Gamma_+}$ and the singular-value form
$\sigma_{\min}/\sigma_{\max}$ of \cref{def:sheafBW} agree numerically
on a randomly generated $A$ of the relevant shape. Two facts are
verified:
\begin{itemize}[leftmargin=*]
\item the linear-algebraic identity
  $\sqrt{\lambda_{\min}(A^* A) / \lambda_{\max}(A^* A)}
  = \sigma_{\min}(A)/\sigma_{\max}(A)$,
  which holds for any $M \times K$ matrix $A$ (a standard consequence
  of the singular value decomposition);
\item the code implementation (\texttt{A\_route\_ch3\_numerical.py})
  computes both forms without introducing numerical artifacts beyond
  IEEE 754 roundoff.
\end{itemize}
What the script does \emph{not} do:
\begin{itemize}[leftmargin=*]
\item It does not provide independent physical evidence that the
  domain $V_?$ in each instance is correctly chosen for the underlying
  application (signal reconstruction in TopSP, finite element solution
  in FEM--EM, band-structure integration in BZ). All three instances
  use the same toy construction
  $U = \mathrm{QR}(\mathrm{randn}(N, K))$, which only matches the
  abstract structural template
  ``finite-dim Hilbert subspace + row-sampled basis representation.''
\item It does not validate the heuristic skeleton of
  \cref{empirical:8b}: that proposition rests on the crossed
  evidence channels of \cref{tab:6-evidence}, not on this script.
\item It does not exercise the cellular-sheaf construction of
  \cref{sec:sheafbw} independently: by \cref{prop:delta-equals-A},
  $\delta^0 = A$ on the sampling sheaf, so
  $(\delta^0)^* \delta^0 = A^* A$ holds by definition rather than as
  a numerical fact. The deviation reported in \cref{tab:7-results}
  is the IEEE 754 roundoff of computing eigenvalues of $A^*A$ versus
  squared singular values of $A$; it does not test the cellular
  sheaf structure.
\end{itemize}
The script is best understood as a sanity check on the algebra of
\cref{def:sheafBW} and on the code path that computes it, not as
empirical support for the framework's applicability to any specific
domain.

% -------------------------------------------------------------
\subsection{Physical evidence per instance}\label{subsec:7-evidence}

The empirical content of each instance rests on the following sources,
which are not exercised by the script of \cref{subsec:7-setup}.

\begin{table}[h]
\centering
\small
\begin{tabular}{@{}l p{0.45\linewidth} p{0.32\linewidth}@{}}
\toprule
\textbf{Instance} & \textbf{Empirical content} & \textbf{Source} \\
\midrule
1: TopSP    & $35$ MRI sampling configurations; full column rank in all $35$;
              Cartesian $\sheafBW \approx 0.0091$ vs.\ radial $\approx 2.07 \times 10^{-5}$
            & this paper, project simulation suite (\cref{subsec:6-topsp}) \\
2: FEM--EM  & finite-dimensional restriction is a project-internal scope
              statement, not an empirical claim
            & \cref{subsec:6-femem} \\
3: BZ       & five decades of condensed-matter practice consistently report
              non-Cartesian tessellations outperforming Cartesian quadrature
            & \cite{Blochl1994,LehmannTaut1972,JepsenAndersen1971,MonkhorstPack1976} \\
\bottomrule
\end{tabular}
\caption{Sources of empirical content per instance. Instance 2 is
  established by the finite-dimensional restriction of
  \cref{subsec:sheafbw-hilbert-complex}; the row labeled ``finite-dimensional
  restriction'' is a project-internal scope statement and is not a
  software- or hardware-derived empirical claim.}
\label{tab:7-evidence}
\end{table}

\noindent
The numerical verification closes the main argument of the paper. We
conclude with a discussion of the two independent open problems left
by this work.


% =============================================================
% Section 8. Discussion and Future Work
% =============================================================
\section{Discussion and Open Problems}\label{sec:discussion}
% =============================================================
% Section 8. Discussion and Open Problems
% Template: paper/v0.1/notes/narrative.md §4 §8
% Structure: §8.1 contributions recap / §8.2 OP 4.5-b /
%            §8.3 OP 8b-rigor / §8.4 MRI engineering outlook /
%            §8.5 acknowledgments
% =============================================================

\noindent
We have constructed a categorical bridge between finite element
exterior calculus, cellular sheaf signal processing, and topological
signal processing, instantiated it in three previously disjoint
domains, and verified its main equality numerically to machine
precision. We now restate the seven contributions in the order they
were established, give the precise content of the two open problems
that remain, and outline the engineering verification on the MRI
side that lies outside this paper.

% -------------------------------------------------------------
\subsection{Summary of contributions}\label{subsec:disc-contrib}

\begin{enumerate}[leftmargin=*,label=(C\arabic*)]
\item The triangle equivalence (\cref{thm:triangle-equiv}) reads one
  scalar---$\sigma_{\min}(A)/\sigma_{\max}(A)$---through three
  schools: FEEC singular value bounds (\cref{subsec:tri-ab}), TopSP
  Riesz frame conditions (\cref{rem:three-readings}), and Robinson
  sheaf cohomology (\cref{subsec:tri-ac}). The proof produces a
  candidate sheaf bandwidth that is not yet categorically intrinsic
  (\cref{def:sheafbw-candidate}), and \cref{subsec:tri-narrowing}
  precisely narrows Robinson's open problem to the construction of an
  intrinsic version.
\item The categorically intrinsic candidate
  $\sheafBW(\mathscr{A}) := \sqrt{\gamma_+(\delta^0)/\Gamma_+(\delta^0)}$
  (\cref{def:sheafBW}) is unitarily invariant
  (\cref{prop:sheafBW-unitary}). The well-definedness of the
  spectral gap quantities relies on the closed Hilbert complex
  structure of \cref{thm:closed-hilbert-complex}, which uses the
  finite-dimensional restriction of \cref{def:finite-dim-sheaf}.
\item On the abstract sampling cell complex $X^A$
  (\cref{def:Xsamp}), the sampling sheaf $\samps$
  (\cref{def:sampling-sheaf}) reduces the intrinsic invariant to the
  ambient ratio $\sigma_{\min}(A)/\sigma_{\max}(A)$ \emph{strictly}
  (\cref{thm:strict-equality}) on the well-conditioned locus
  $\sigma_{\min}(A) > 0$. The double $0$-cell construction
  $\{\sigma_B, \sigma_\infty\}$ is what places $X^A$ into the
  Hansen--Ghrist regularity framework
  \cite[Def.~1, Cond.~4]{HansenGhrist2019};
  \cref{rem:rank-deficient-locus} delineates the rank-deficient
  boundary on which the strict equality fails.
\item Three previously disjoint domains
  (\cref{sec:three_instances}) fit a common categorical template:
  TopSP ($V_B^{\TopSP}$, \cref{subsec:6-topsp}), FEM--EM ($V_S^{\FEM}$,
  \cref{subsec:6-femem}), and Brillouin zone tetrahedron quadrature
  ($V_{\occ}^{\BZ}$, \cref{subsec:6-bz}). The three-instance
  comparison is recorded in \cref{tab:6-instances}.
\item The empirical proposition (\cref{empirical:8b}) formulates the
  $S$--$I$ Pareto trade-off as a directional empirical statement
  supported by two empirical channels and a heuristic skeleton
  (\cref{tab:6-evidence}); it is not promoted to theorem status, and
  it is not invoked in the proof of \cref{thm:strict-equality}---the
  categorical bridge is independent of it.
\item The numerical verification (\cref{sec:numerical}) records
  $\mathrm{reldiff} \le 1.64 \times 10^{-14}$ across the three toy
  instances at $\mathtt{seed} = 20260507$ (with Instances~1 and~3
  attaining $\le 8 \times 10^{-16}$), with explicit scope: the script
  verifies the SVD-based identity and the code path, not the
  per-instance physical evidence
  (\cref{subsec:7-honest-statement,tab:7-evidence}).
\item Two independent open problems remain, formulated below.
\end{enumerate}

% -------------------------------------------------------------
\subsection{Open Problem: non-trivial dagger on the full category}\label{subsec:disc-op-dagger}

\Cref{op:non-trivial-dagger} asks whether the unitary subcategory
groupoid of \cref{def:unitary-subcat} extends to a non-trivial
$\dagger$ structure on the full category $\CellSh^{\fin}_{\Hilb}(X)$.
The substantive content of this question is the tension between two
candidate lifts: dagger compactness in $\Hilb$
\cite[\S 2.3]{HeunenVicary2019} does not lift directly to cellular
sheaves because the restriction maps are not in general isometries,
while the dagger mono structure of
\cite[\S 3]{LahtinenStenvall2020} does lift but produces only the
unitary subgroupoid rather than a structure on the full category.
\Cref{rem:trivial-dagger} explains why the naive lift collapses; the
construction of a non-trivial extension is left open.
\Cref{op:non-trivial-dagger} is independent of all theorems of this
paper.

% -------------------------------------------------------------
\subsection{Open Problem: constructive Pareto proof}\label{subsec:disc-op-pareto}

\Cref{op:8b-rigor} asks for a constructive Pareto-optimality proof
upgrading the empirical proposition \cref{empirical:8b} to a
theorem. The proposed mode of attack proceeds in three steps:
compactness of the $V$-parameter space under suitable boundary
conditions; convex relaxation of the joint map
$V \mapsto (S(X^A_V), I(X^A_V))$, with the dual optimization
encoding the trade-off; and closing the duality gap via spectral
arguments paralleling \cite{Pilleboue2015}. Resolution of this
problem would replace \cref{empirical:8b} by a theorem and the
heuristic skeleton by a constructive proof. \Cref{op:8b-rigor} is
independent of the categorical bridge and of all theorems of this
paper.

% -------------------------------------------------------------
\subsection{Engineering outlook on the MRI side}\label{subsec:disc-mri-outlook}

The triangle equivalence (\cref{thm:triangle-equiv}) and the strict
equality (\cref{thm:strict-equality}) reduce, on the MRI sampling
problem, to a singular value ratio computable directly from the
acquisition matrix $A$. The project-internal numerical study
recorded in \cref{tab:6-evidence} examines $35$ sampling
configurations at $N = 200$; the multi-$N$ scan
($N \in \{200, 300, 1000\}$) shows the top three configurations
(Cartesian, random, sparkling) and the bottom two (spiral, radial)
separated by two to three orders of magnitude in the singular value
ratio across the entire range.
A larger-scale engineering verification---with $N \approx 10^4$
acquisition points, real phantom data, and one representative
baseline drawn from each of the four schools surveyed in
\cref{sec:related}---is the natural next step on the MRI side. That
engineering verification is outside the scope of this paper.

% -------------------------------------------------------------
\subsection*{Acknowledgments}

[Acknowledgments placeholder.]


% =============================================================
% Bibliography
% =============================================================
\bibliographystyle{amsalpha}
\bibliography{bib/refs}

\end{document}
